% Section file for Chapter: Twisted Holography and Quantum Gravity
% Result (G6): Subleading Soft Graviton Theorems from the Shadow Tower
%
% Dependencies:
%   - thm:mc2-bar-intrinsic (bar-intrinsic MC element)
%   - def:shadow-postnikov-tower (shadow Postnikov tower)
%   - def:modular-shadow-connection (modular shadow connection)
%   - prop:master-equation-from-mc (all-arity master equation)
%   - thm:shadow-connection-kz (KZ from shadow connection)
%   - thm:cubic-gauge-triviality (cubic gauge triviality)
%   - thm:w-virasoro-quintic-forced (Virasoro quintic forced)
%
% Status: \ClaimStatusProvedHere for the algebraic theorems;
%   celestial holography interpretation is physical/structural.

\section{Soft graviton theorems from the shadow tower}
\label{sec:thqg-soft-graviton-theorems}
\index{soft graviton theorem|textbf}
\index{Ward identity!gravitational soft theorem}
\index{shadow Postnikov tower!soft graviton theorems}
\index{celestial holography!from modular shadow connection}

% ======================================================================
% LOCAL MACROS (providecommand: no conflict with main.tex preamble)
% ======================================================================
\providecommand{\Sh}{\operatorname{Sh}}
\providecommand{\cA}{\mathcal{A}}
\providecommand{\cV}{\mathcal{V}}
\providecommand{\MC}{\text{MC}}
\providecommand{\gAmod}{\mathfrak{g}_{\cA}^{\mathrm{mod}}}
\providecommand{\Defcyc}{\operatorname{Def}_{\mathrm{cyc}}}
\providecommand{\dzero}{d_0}
\providecommand{\barBch}{\bar{B}^{\text{ch}}}
\providecommand{\Omegach}{\Omega^{\text{ch}}}
\providecommand{\dfib}{d_{\mathrm{fib}}}
\providecommand{\Ydg}{\mathcal{Y}^{\mathrm{dg}}}
\providecommand{\Gmod}{\mathcal{G}^{\mathrm{mod}}}
\providecommand{\Convstr}{\operatorname{Conv}_{\mathrm{str}}}
\providecommand{\Convinf}{\operatorname{Conv}_{\infty}}
\providecommand{\Definfmod}{\operatorname{Def}_{\infty}^{\mathrm{mod}}}
\providecommand{\Res}{\operatorname{Res}}
\providecommand{\Tr}{\operatorname{Tr}}
\providecommand{\ad}{\operatorname{ad}}
\providecommand{\id}{\mathrm{id}}
\providecommand{\rank}{\operatorname{rank}}
\providecommand{\wt}{\operatorname{wt}}
\providecommand{\Hom}{\operatorname{Hom}}
\providecommand{\End}{\operatorname{End}}
\providecommand{\Sym}{\operatorname{Sym}}
\providecommand{\Aut}{\operatorname{Aut}}
\providecommand{\sgn}{\operatorname{sgn}}
\providecommand{\Dg}[1]{D_{#1}}
\providecommand{\mcurv}[1]{m_0^{(#1)}}
\providecommand{\colldiv}[2]{D_{#1#2}}
\providecommand{\LogForm}[2]{\eta_{#1#2}}
\providecommand{\fg}{\mathfrak{g}}
\providecommand{\cW}{\mathcal{W}}

% ======================================================================
%
%   SECTION OVERVIEW
%
% The modular shadow connection
%   nabla^hol_{g,n} = d - Sh_{g,n}(Theta_A)
% is a flat connection on the space of derived conformal blocks
% V_{g,n}(A).  Its flatness is a consequence of the MC equation
% for Theta_A (Theorem thm:mc2-bar-intrinsic).  When the arity
% decomposition of Sh_{g,n} is expanded, each arity contributes
% a Ward identity.  These are the soft graviton theorems.
%
%   Arity 2: Weinberg leading soft theorem
%   Arity 3: Cachazo--Strominger subleading
%   Arity 4: sub-subleading from quartic contact Q^contact
%   Arity r: all-arity nonlinear soft theorem from Sh_r
%
% The shadow depth r_max(A) controls how many independent soft
% theorems exist.  The infinite tower for Virasoro/W_N algebras
% is the gravitational instance.
%
% ======================================================================

The central formula of the holographic modular Koszul datum
(Definition~\ref{def:holographic-modular-koszul-datum}) is the modular
shadow connection
\begin{equation}\label{eq:thqg-VI-shadow-connection-master}
\nabla^{\mathrm{hol}}_{g,n}
\;=\;
d \;-\; \Sh_{g,n}(\Theta_\cA),
\end{equation}
a flat connection on the space of derived conformal blocks
$\cV_{g,n}(\cA)$
(Definition~\ref{def:modular-shadow-connection}).
Flatness
\begin{equation}\label{eq:thqg-VI-flatness}
\bigl(\nabla^{\mathrm{hol}}_{g,n}\bigr)^2 = 0
\end{equation}
is the $(g,n)$-projection of the Maurer--Cartan equation
$D_\cA\Theta_\cA + \tfrac{1}{2}[\Theta_\cA,\Theta_\cA] = 0$
(Theorem~\ref{thm:mc2-bar-intrinsic}).  The horizontal sections of
$\nabla^{\mathrm{hol}}_{g,n}$ are the genus-$g$, $n$-point
correlators of the holographic theory.  The Ward identities of this
connection, decomposed by the arity filtration of the shadow Postnikov
tower (Definition~\ref{def:shadow-postnikov-tower}), are the soft
graviton theorems of the title.

The purpose of this section is fourfold:
\begin{enumerate}[label=\textup{(\arabic*)}]
\item to prove that the arity-$2$ Ward identity of
  $\nabla^{\mathrm{hol}}_{g,n}$ reproduces the Weinberg leading soft
  graviton theorem and its BMS extension
  (\S\ref{subsec:thqg-VI-leading-soft});
\item to prove that the arity-$3$ Ward identity gives the
  Cachazo--Strominger subleading soft graviton theorem, with the cubic
  gauge triviality of Theorem~\ref{thm:cubic-gauge-triviality}
  controlling when the subleading theorem is non-trivial
  (\S\ref{subsec:thqg-VI-subleading-soft});
\item to prove that the arity-$4$ Ward identity produces the
  sub-subleading soft graviton theorem with explicit coefficient
  $Q^{\mathrm{contact}}_{\mathrm{Vir}} = 10/[c(5c+22)]$, and to
  derive the clutching law that controls its factorization across
  degenerate curves
  (\S\ref{subsec:thqg-VI-sub-subleading-soft});
\item to establish the infinite tower of higher soft theorems for
  algebras with $r_{\max}(\cA) = \infty$ (Virasoro and $\mathcal{W}_N$
  families), prove $o_5(\mathrm{Vir}_c) \neq 0$, and develop the
  celestial holography interpretation
  (\S\ref{subsec:thqg-VI-higher-soft},
  \S\ref{subsec:thqg-VI-celestial}).
\end{enumerate}

Throughout this section, $\cA$ is a modular Koszul chiral algebra on a
smooth projective curve~$X$, and $\Theta_\cA \in \MC(\gAmod)$ is the
universal MC element of
Theorem~\ref{thm:mc2-bar-intrinsic}.  The shadow depth is
$r_{\max}(\cA) := \sup\{r \geq 2 : \cA^{\mathrm{sh}}_{r,0} \neq 0\}$,
and the shadow Postnikov tower is
$\Theta_\cA^{\leq 2} \to \Theta_\cA^{\leq 3} \to
\Theta_\cA^{\leq 4} \to \cdots \to \Theta_\cA$.
All gradings are cohomological ($|d| = +1$).

% ======================================================================
%
%   SUBSECTION 1: SHADOW CONNECTION AND ARITY DECOMPOSITION
%
% ======================================================================

\subsection{The shadow connection and its arity decomposition}
\label{subsec:thqg-VI-shadow-connection}
\index{shadow connection!arity decomposition}
\index{arity filtration!shadow connection}

\subsubsection{The shadow representation}

\begin{definition}[Shadow representation at $(g,n)$]
\label{def:thqg-VI-shadow-representation}
\index{shadow representation|textbf}
For each pair $(g,n)$ with $2g - 2 + n > 0$, the
\emph{shadow representation at $(g,n)$} is the linear map
\begin{equation}\label{eq:thqg-VI-shadow-representation}
\Sh_{g,n} \colon
\Defcyc^{\mathrm{mod}}(\cA)
\;\longrightarrow\;
\End\bigl(\cV_{g,n}(\cA)\bigr)
\end{equation}
defined as follows.  An element
$\varphi \in \Defcyc^{\mathrm{mod}}(\cA)$
of genus~$h$ and arity~$r$ acts on a conformal block
$\Phi \in \cV_{g,n}(\cA; V_1, \dotsc, V_n)$ by evaluating the
multilinear operation $\varphi$ on~$r$ insertions drawn from the
$n$-point configuration, paired with the moduli-space coefficient
from $C_\bullet(\overline{\mathcal{M}}_{g,n+1})$, and summed over
all ways to select the $r$~insertions from the $n$~marked points.

Concretely, if $\varphi$ has arity~$r$ and genus~$h$:
\begin{equation}\label{eq:thqg-VI-shadow-action-concrete}
\bigl(\Sh_{g,n}(\varphi) \cdot \Phi\bigr)(z_1,\dotsc,z_n)
\;=\;
\sum_{\substack{I \subseteq \{1,\dotsc,n\} \\ |I| = r}}
\int_{\overline{\mathcal{M}}_{h,r+1}}
\varphi\bigl(\phi_{i_1}(z_{i_1}),\dotsc,\phi_{i_r}(z_{i_r})\bigr)
\;\Phi\bigl|_{I \to \varphi(I)},
\end{equation}
where the notation $\Phi|_{I \to \varphi(I)}$ denotes the conformal
block with the $r$~insertions in~$I$ replaced by the output
of~$\varphi$, and the moduli integral produces the coefficient-valued
endomorphism.
\end{definition}

\begin{remark}[Arity filtration on the connection form]
\label{rem:thqg-VI-arity-filtration}
\index{shadow connection!filtration by arity}
The modular convolution algebra $\gAmod$ carries a natural
\emph{arity filtration}: $F^{\leq r}\gAmod$ consists of elements
$\varphi$ with arity at most~$r$.  Since $\Theta_\cA$ decomposes
as $\Theta_\cA = \sum_{r \geq 2} \Theta_\cA^{(r)}$ where
$\Theta_\cA^{(r)} \in F^r\gAmod / F^{r-1}\gAmod$ is the arity-$r$
component, the shadow connection form decomposes as
\begin{equation}\label{eq:thqg-VI-arity-decomposition}
\Sh_{g,n}(\Theta_\cA)
\;=\;
\sum_{r=2}^{\infty}
\Sh_{g,n}\bigl(\Theta_\cA^{(r)}\bigr)
\;=\;
\Sh_{g,n}\bigl(\kappa(\cA)\bigr)
\;+\;
\Sh_{g,n}\bigl(\mathfrak{C}(\cA)\bigr)
\;+\;
\Sh_{g,n}\bigl(\mathfrak{Q}(\cA)\bigr)
\;+\;
\sum_{r=5}^{\infty}\Sh_{g,n}\bigl(\Sh_r(\cA)\bigr).
\end{equation}
The notation uses:
$\Theta_\cA^{(2)} = \kappa(\cA)$ (the modular characteristic,
Theorem~D),
$\Theta_\cA^{(3)} = \mathfrak{C}(\cA)$ (the cubic shadow),
$\Theta_\cA^{(4)} = \mathfrak{Q}(\cA)$ (the quartic shadow /
contact invariant), and
$\Theta_\cA^{(r)} = \Sh_r(\cA) = \cA^{\mathrm{sh}}_{r,0}$
for general~$r$.

For an algebra of shadow depth $r_{\max}(\cA) < \infty$,
the sum~\eqref{eq:thqg-VI-arity-decomposition} terminates at
$r = r_{\max}$.

The \emph{shadow depth classes} are defined in
\S\ref{sec:thqg-gravitational-complexity}
(Definition~\ref{def:thqg-shadow-archetype}):
\textbf{G}~(Gaussian, $r_{\max}=2$),
\textbf{L}~(Lie/tree, $r_{\max}=3$),
\textbf{C}~(Contact, $r_{\max}=4$),
\textbf{M}~(Mixed, $r_{\max}=\infty$).
The class of $\cA$ determines how many terms in the
arity decomposition are nonzero.
\end{remark}

\begin{proposition}[Arity decomposition of the shadow connection;
\ClaimStatusProvedHere]
\label{prop:thqg-VI-arity-decomposition}
\index{shadow connection!arity decomposition|textbf}
The shadow connection $\nabla^{\mathrm{hol}}_{g,n}$ admits a canonical
decomposition
\begin{equation}\label{eq:thqg-VI-nabla-arity}
\nabla^{\mathrm{hol}}_{g,n}
\;=\;
\nabla^{(2)}_{g,n}
\;+\;
A^{(3)}_{g,n}
\;+\;
A^{(4)}_{g,n}
\;+\;
\sum_{r \geq 5} A^{(r)}_{g,n},
\end{equation}
where:
\begin{enumerate}[label=\textup{(\roman*)}]
\item $\nabla^{(2)}_{g,n} := d - \Sh_{g,n}(\kappa(\cA))$ is
  the \emph{leading shadow connection}, a flat connection on
  $\cV_{g,n}(\cA)$ that depends only on the modular
  characteristic~$\kappa(\cA)$;
\item $A^{(r)}_{g,n} := -\Sh_{g,n}(\Theta_\cA^{(r)})$ is
  the \emph{arity-$r$ connection perturbation} for $r \geq 3$,
  an $\End(\cV_{g,n})$-valued $1$-form on the moduli of
  configurations.
\end{enumerate}
\end{proposition}

\begin{proof}
This is immediate from the decomposition
$\Theta_\cA = \sum_{r \geq 2} \Theta_\cA^{(r)}$
and the linearity of the shadow representation
$\Sh_{g,n}$ (Definition~\ref{def:thqg-VI-shadow-representation}).
The assertion that $\nabla^{(2)}_{g,n}$ is itself flat when
$\cA$ is Gaussian (class~$\mathbf{G}$, i.e.~$r_{\max} = 2$) follows
from the MC equation restricted to arity~$2$: the quadratic piece
$\tfrac{1}{2}[\kappa,\kappa]$ vanishes because $\kappa$ is scalar
(Theorem~\ref{thm:genus-universality}).
\end{proof}

\subsubsection{The flatness equation by arity}

\begin{theorem}[Flatness by arity; \ClaimStatusProvedHere]
\label{thm:thqg-VI-flatness-by-arity}
\index{flatness!arity decomposition|textbf}
\index{shadow connection!flatness by arity}
The flatness condition
$(\nabla^{\mathrm{hol}}_{g,n})^2 = 0$
decomposes, under the arity filtration, into the following
system of equations.  For each $r \geq 2$, define the
\emph{arity-$r$ flatness residual}
\begin{equation}\label{eq:thqg-VI-flatness-arity-r}
F^{(r)}_{g,n}
\;:=\;
\Bigl(
(\nabla^{\mathrm{hol}}_{g,n})^2
\Bigr)_{r}
\;=\; 0,
\end{equation}
where $(\cdot)_r$ denotes projection to the arity-$r$ component
of $\End(\cV_{g,n})$.  Then:
\begin{enumerate}[label=\textup{(\alph*)}]
\item \emph{Arity~$2$}:
\begin{equation}\label{eq:thqg-VI-flat-arity-2}
\nabla^{(2)}_{g,n} \circ \nabla^{(2)}_{g,n} = 0
\qquad\Longleftrightarrow\qquad
d\bigl(\Sh_{g,n}(\kappa)\bigr) = 0,
\quad
\Sh_{g,n}(\kappa) \wedge \Sh_{g,n}(\kappa) = 0.
\end{equation}
This is the \emph{leading soft constraint}: the modular
characteristic acts by commuting endomorphisms.
\item \emph{Arity~$3$}:
\begin{equation}\label{eq:thqg-VI-flat-arity-3}
\nabla^{(2)}_{g,n}\bigl(A^{(3)}_{g,n}\bigr)
\;+\;
o_3(\cA)_{g,n}
\;=\; 0,
\end{equation}
where $o_3(\cA)_{g,n} = \Sh_{g,n}(o_3(\cA))$ is the genus-$(g,n)$
image of the cubic obstruction class.  This is the
\emph{subleading soft constraint}.
\item \emph{Arity~$4$}:
\begin{equation}\label{eq:thqg-VI-flat-arity-4}
\nabla^{(2)}_{g,n}\bigl(A^{(4)}_{g,n}\bigr)
\;+\;
\bigl[A^{(3)}_{g,n},\, A^{(3)}_{g,n}\bigr]
\;+\;
o_4(\cA)_{g,n}
\;=\; 0.
\end{equation}
The quadratic term $[A^{(3)}, A^{(3)}]$ is the self-interaction
of the cubic shadow.  This is the \emph{sub-subleading soft
constraint}.
\item \emph{General arity~$r$}:
\begin{equation}\label{eq:thqg-VI-flat-arity-r-general}
\nabla^{(2)}_{g,n}\bigl(A^{(r)}_{g,n}\bigr)
\;+\;
\sum_{\substack{s+t = r \\ s,t \geq 3}}
\bigl[A^{(s)}_{g,n},\, A^{(t)}_{g,n}\bigr]
\;+\;
o_r(\cA)_{g,n}
\;=\; 0.
\end{equation}
This is the all-arity master equation
(Proposition~\ref{prop:master-equation-from-mc}), represented on
conformal blocks.
\end{enumerate}
\end{theorem}

\begin{proof}
Expand $(\nabla^{\mathrm{hol}}_{g,n})^2$ using the
arity decomposition~\eqref{eq:thqg-VI-nabla-arity}.  Write
$\omega = \sum_{r \geq 2} \Sh_{g,n}(\Theta^{(r)}_\cA)$.
Then
\[
(\nabla^{\mathrm{hol}}_{g,n})^2
= (d - \omega)^2
= -d\omega + \omega \wedge \omega
= -\sum_{r \geq 2} d\bigl(\Sh_{g,n}(\Theta^{(r)})\bigr)
  + \sum_{r,s \geq 2}
    \Sh_{g,n}(\Theta^{(r)}) \wedge \Sh_{g,n}(\Theta^{(s)}).
\]
The total vanishing is the MC equation for $\Theta_\cA$.
Projecting to the arity-$r$ component, the term $d(\Sh_{g,n}(\Theta^{(r)}))$
contributes the covariant derivative $\nabla^{(2)}(A^{(r)})$ (for
$r \geq 3$) plus a correction from the commutator of the leading
connection with $A^{(r)}$.  The quadratic terms
$\Sh_{g,n}(\Theta^{(s)}) \wedge \Sh_{g,n}(\Theta^{(t)})$ with
$s + t = r$ give the bracket terms.  The remainder at arity~$r$
is exactly the obstruction class $o_r(\cA)$ evaluated on conformal
blocks.

More precisely, the all-arity master equation
$\nabla_H(\Sh_r) + \mathfrak{o}^{(r)} = 0$
(Proposition~\ref{prop:master-equation-from-mc}) is obtained by
projecting $D_\cA\Theta_\cA + \tfrac{1}{2}[\Theta_\cA,\Theta_\cA] = 0$
onto the arity-$r$ component of $\Defcyc^{\mathrm{mod}}(\cA)$.
Applying the shadow representation $\Sh_{g,n}$ converts this into
the endomorphism equation~\eqref{eq:thqg-VI-flat-arity-r-general}
on $\cV_{g,n}(\cA)$.  The identification of the
$\nabla^{(2)}$-covariant derivative with the Hodge-linear
part~$\nabla_H$ is the definition of the leading connection.
\end{proof}

\begin{remark}[Arity truncation and the depth of soft gravity]
\label{rem:thqg-VI-arity-truncation}
\index{shadow depth!soft graviton hierarchy}
For an algebra of shadow depth $r_{\max}(\cA)$, the system
\eqref{eq:thqg-VI-flat-arity-r-general} closes at $r = r_{\max}$:
$A^{(r)}_{g,n} = 0$ for all $r > r_{\max}$.  The number of
independent soft graviton theorems is therefore $r_{\max}(\cA) - 1$:
\begin{center}
\small
\renewcommand{\arraystretch}{1.15}
\begin{tabular}{lccc}
\toprule
\emph{Family} &
$r_{\max}$ &
\emph{\# soft theorems} &
\emph{Class} \\
\midrule
Heisenberg / lattice &
  $2$ & $1$ (leading only) & $\mathbf{G}$ \\
Affine $\widehat{\mathfrak{g}}_k$ &
  $3$ & $2$ (leading + subleading) & $\mathbf{L}$ \\
$\beta\gamma$ system &
  $4$ & $3$ (leading + sub + sub-sub) & $\mathbf{C}$ \\
Virasoro / $\mathcal{W}_N$ &
  $\infty$ & $\infty$ & $\mathbf{M}$ \\
\bottomrule
\end{tabular}
\end{center}
This is the gravitational interpretation of shadow complexity
(Result~\textbf{(G2)}, \S\ref{sec:thqg-gravitational-complexity}).
\end{remark}

\subsubsection{Ward identities from the shadow connection}

\begin{definition}[Soft shadow Ward identity at arity~$r$]
\label{def:thqg-VI-soft-ward-identity}
\index{soft Ward identity|textbf}
\index{Ward identity!soft shadow}
Let $\cA$ be a modular Koszul chiral algebra and let
$\Phi \in \cV_{g,n}(\cA; V_1, \dotsc, V_n)$ be a horizontal
section of $\nabla^{\mathrm{hol}}_{g,n}$, i.e.\ a correlator:
$\nabla^{\mathrm{hol}}_{g,n}\,\Phi = 0$.
The \emph{soft shadow Ward identity at arity~$r$} is the
constraint on~$\Phi$ obtained by projecting the horizontality
condition to the arity-$r$ component:
\begin{equation}\label{eq:thqg-VI-ward-arity-r}
\boxed{%
\bigl(\Sh_{g,n}(\Theta^{(r)}_\cA)\bigr) \cdot \Phi
\;=\;
\text{(contributions from arities $< r$)}.}
\end{equation}
When $r = 2$, the right side vanishes and the identity is the
\emph{leading soft Ward identity}.  For $r \geq 3$, the right side
involves the lower-arity shadows acting on~$\Phi$ through the
bracket terms of the flatness equation.
\end{definition}

\begin{theorem}[Shadow Ward identities determine correlators;
\ClaimStatusProvedHere]
\label{thm:thqg-VI-ward-determines-correlators}
\index{shadow Ward identity!determination of correlators|textbf}
Let $\cA$ be a modular Koszul chiral algebra with shadow depth
$r_{\max}(\cA) \leq \infty$.  The system of shadow Ward identities
at arities $r = 2, 3, \dotsc, r_{\max}$ is equivalent to the
horizontality condition $\nabla^{\mathrm{hol}}_{g,n}\,\Phi = 0$.
In particular:
\begin{enumerate}[label=\textup{(\roman*)}]
\item If $r_{\max} < \infty$, the system is a \emph{finite}
  system of differential equations on genus-$g$, $n$-point
  correlators, determined by the finite shadow tower
  $\{\Sh_2, \dotsc, \Sh_{r_{\max}}\}$.
\item If $r_{\max} = \infty$, the system is an infinite tower
  of coupled differential equations whose consistency is
  guaranteed by the MC equation.
\item For $\cA = \widehat{\mathfrak{g}}_k$:
  the leading Ward identity recovers the KZ equation, and the
  subleading identity is automatically satisfied by the Jacobi
  identity of~$\mathfrak{g}$
  (Theorem~\ref{thm:shadow-connection-kz}).
\item For $\cA = \mathrm{Vir}_c$: the full infinite tower of
  Ward identities encodes the BPZ conformal Ward identities
  together with all their nonlinear extensions.
\end{enumerate}
\end{theorem}

\begin{proof}
The horizontality condition $\nabla^{\mathrm{hol}}_{g,n}\Phi = 0$
is the single equation $(d - \Sh_{g,n}(\Theta_\cA))\Phi = 0$.
Since $\Theta_\cA = \sum_{r \geq 2} \Theta_\cA^{(r)}$ and the
shadow representation is additive, the arity-$r$ projection of
$\Sh_{g,n}(\Theta_\cA)\Phi$ extracts exactly the arity-$r$ Ward
identity.  The full horizontality condition is the conjunction of
these projections for all~$r$.

For~(i): when $r_{\max} < \infty$, only finitely many arities
contribute and the system is a finite differential system.

For~(ii): when $r_{\max} = \infty$, the consistency of the
infinite system follows from the flatness
$(\nabla^{\mathrm{hol}})^2 = 0$, which is the MC equation.

For~(iii): the arity-$2$ shadow of $\Theta_{\widehat{\fg}_k}$ is
$\kappa(\widehat{\fg}_k) = k\,\delta_{ab}/2$ (the Killing form at
level~$k$).  The shadow representation at genus~$0$ acts by
\begin{equation}\label{eq:thqg-VI-kz-recovery}
\Sh_{0,n}(\kappa) \cdot \Phi
= \sum_{i < j} \frac{t^a_i\,t^a_j}{z_i - z_j}\;\Phi,
\end{equation}
where $t^a_i$ denotes the action of $\mathfrak{g}$ on the
$i$-th module.  This is the KZ connection form.  The arity-$3$
Ward identity involves $\mathfrak{C}(\widehat{\fg}_k)$, which
is the cubic shadow.  For the affine algebra, $o_4 = 0$ by the
Jacobi identity
(Example~\ref{ex:tower-four-archetypes}): the quartic obstruction
vanishes because the structure constants satisfy
$f^a_{bc}\,f^c_{de} + \text{cyc} = 0$.  Hence the tower terminates
at $r_{\max} = 3$ and the subleading identity is the last one.

For~(iv): The Virasoro algebra has $r_{\max} = \infty$
(Theorem~\ref{thm:w-virasoro-quintic-forced}), so the tower
never terminates.  The arity-$2$ Ward identity at genus~$0$
is the conformal Ward identity involving
$\kappa(\mathrm{Vir}_c) = c/2$,
which gives the standard $T(z)$ OPE.  The higher arities encode the
nonlinear constraints arising from normal-ordered products
$:T^n:$ at all orders.
\end{proof}

\begin{remark}[Comparison with the Weinberg--Cachazo--Strominger
hierarchy]
\label{rem:thqg-VI-wcs-comparison}
\index{Weinberg soft theorem!shadow connection interpretation}
\index{Cachazo--Strominger!subleading soft theorem}
In the gravitational scattering amplitude literature, the soft
graviton theorems form a hierarchy:
\begin{enumerate}[label=\textup{(\arabic*)}]
\item \emph{Leading} (Weinberg 1965): the $\omega^{-1}$ pole of
  the soft factor is universal, determined by total energy-momentum.
\item \emph{Subleading} (Cachazo--Strominger 2014): the $\omega^0$
  term is determined by angular momentum.
\item \emph{Sub-subleading} (Cachazo--He--Yuan 2014): the
  $\omega^1$ term involves the stress tensor.
\item Higher: (Hamada--Shiu 2018, Li--Strominger 2018): the
  $\omega^k$ terms are progressively less universal.
\end{enumerate}
The shadow-tower decomposition gives a clean algebraic explanation:
the arity-$r$ shadow at genus~$0$ produces the order-$(r-2)$
term in the soft expansion.  The modular characteristic $\kappa$
(arity~$2$) gives the leading $\omega^{-1}$ term; the cubic shadow
$\mathfrak{C}$ (arity~$3$) gives the subleading $\omega^0$ term;
the quartic contact $\mathfrak{Q}$ (arity~$4$) gives the
sub-subleading $\omega^1$ term.  The shadow depth $r_{\max}$
is the algebraic counterpart of the maximal non-trivial order in
the soft expansion.  See Table~\ref{tab:thqg-VI-soft-dictionary}.
\end{remark}

\begin{remark}[Genus dependence of the Ward identities]
\label{rem:thqg-VI-genus-dependence}
\index{soft Ward identity!genus dependence}
The genus-$g$ Ward identities at arity~$r$ involve the full
genus-$g$ shadow $\Sh_{g,n}(\Theta_\cA^{(r)})$, which depends on
both the algebraic shadow $\cA^{\mathrm{sh}}_{r,0}$ and the
moduli-space integral over $\overline{\mathcal{M}}_{g,n+1}$.
At genus~$0$, the moduli integral localizes to the
Fulton--MacPherson compactification of configuration space, and the
Ward identities take the familiar form of differential equations
in the insertion points~$z_i$.
At genus~$g \geq 1$, the Ward identities involve period integrals
over the moduli space and are coupled to the modular shadow
connection on $\overline{\mathcal{M}}_g$ itself.  The genus
spectral sequence
(Construction~\ref{const:vol1-genus-spectral-sequence})
organizes this coupling.
\end{remark}

% ======================================================================
%
%   SUBSECTION 2: LEADING SOFT THEOREM FROM KAPPA
%
% ======================================================================

\subsection{The leading soft theorem from the modular characteristic}
\label{subsec:thqg-VI-leading-soft}
\index{leading soft graviton theorem|textbf}
\index{Weinberg soft theorem!from kappa}

The leading soft graviton theorem is the statement that the
emission of a soft graviton from an $n$-point amplitude is
controlled by a universal factor depending only on the
energy-momentum of the hard particles.  In the shadow-tower
framework, this is the arity-$2$ Ward identity of the modular
shadow connection, controlled by the modular
characteristic~$\kappa(\cA)$.

\subsubsection{The modular characteristic and its action on
conformal blocks}

\begin{proposition}[Action of $\kappa$ on conformal blocks;
\ClaimStatusProvedHere]
\label{prop:thqg-VI-kappa-action}
\index{modular characteristic!action on conformal blocks|textbf}
Let $\cA$ be a modular Koszul chiral algebra with modular
characteristic $\kappa(\cA) \in \cA^{\mathrm{sh}}_{2,0}$.
At genus~$0$, the shadow representation of $\kappa(\cA)$ on
$n$-point conformal blocks $\cV_{0,n}(\cA; V_1, \dotsc, V_n)$
is given by
\begin{equation}\label{eq:thqg-VI-kappa-genus0}
\Sh_{0,n}(\kappa) \cdot \Phi(z_1, \dotsc, z_n)
\;=\;
\sum_{1 \leq i < j \leq n}
\frac{\kappa(\phi_i, \phi_j)}{z_i - z_j}\;
\Phi(z_1, \dotsc, z_n),
\end{equation}
where $\kappa(\phi_i, \phi_j)$ denotes the bilinear pairing
determined by the binary shadow $\kappa \in \Hom_{S_2}(V \otimes V, \mathbb{C})$,
acting on the $i$-th and $j$-th insertions.

At genus~$g \geq 1$, the shadow representation acquires a period
correction:
\begin{equation}\label{eq:thqg-VI-kappa-genus-g}
\Sh_{g,n}(\kappa) \cdot \Phi
\;=\;
\biggl(\sum_{i < j}
\frac{\kappa(\phi_i, \phi_j)}{z_i - z_j}
\;+\;
\kappa(\cA)\;\omega_g\biggr)\Phi,
\end{equation}
where $\omega_g \in H^0(\overline{\mathcal{M}}_g, \mathbb{E}^\vee)$
is the Hodge class and the second term is the curvature correction
$\dfib^2 = \kappa(\cA) \cdot \omega_g$ from the genus-$g$ bar
complex (Theorem~\ref{thm:modular-characteristic}).
\end{proposition}

\begin{proof}
The genus-$0$ formula~\eqref{eq:thqg-VI-kappa-genus0} follows
from the definition of the shadow representation
(Definition~\ref{def:thqg-VI-shadow-representation}): at arity~$2$,
the sum is over all pairs $(i,j)$, and the moduli integral over
$\overline{\mathcal{M}}_{0,3} \cong \{\mathrm{pt}\}$ is trivial.
The propagator on $\overline{C}_2(\mathbb{P}^1) = \mathbb{P}^1$
gives the factor $1/(z_i - z_j)$.

At genus~$g \geq 1$, the moduli integral over
$\overline{\mathcal{M}}_{g,3}$ (the universal curve with one
additional marked point) splits into a configuration-space piece
(giving the same sum over pairs) and a modular piece
(giving the Hodge class~$\omega_g$).  The latter is precisely
the curvature of the genus-$g$ bar complex:
$\dfib^2 = \kappa(\cA) \cdot \omega_g$.
\end{proof}

\subsubsection{The Weinberg soft theorem}

\begin{theorem}[Leading soft graviton theorem; \ClaimStatusProvedHere]
\label{thm:thqg-VI-leading-soft}
\index{leading soft graviton theorem!proof|textbf}
\index{Weinberg soft theorem!modular proof}
Let $\cA$ be a modular Koszul chiral algebra, and let
$\Phi \in \cV_{0,n+1}(\cA; V_1, \dotsc, V_n, V_s)$ be an
$(n+1)$-point genus-$0$ correlator with the last insertion
$\phi_s(z_s)$ taken to be a soft insertion: a state in the
arity-$2$ sector of the shadow algebra, normalized so that
$\kappa(\phi_s, \cdot) = \id_{V_i}$ for each hard insertion
$\phi_i$.  Then the leading soft limit is
\begin{equation}\label{eq:thqg-VI-weinberg}
\lim_{z_s \to \infty}
z_s \cdot \Phi(z_1, \dotsc, z_n, z_s)
\;=\;
S^{(0)}_{\mathrm{Weinberg}}
\;\cdot\;
\Phi_n(z_1, \dotsc, z_n),
\end{equation}
where $\Phi_n$ is the $n$-point correlator and the
\emph{Weinberg soft factor} is
\begin{equation}\label{eq:thqg-VI-weinberg-factor}
S^{(0)}_{\mathrm{Weinberg}}
\;=\;
\sum_{i=1}^{n}
\frac{\kappa(\phi_s, \phi_i)}{z_s - z_i}
\;=\;
\sum_{i=1}^{n}
\frac{p_i \cdot \varepsilon_s}{p_i \cdot q_s},
\end{equation}
with the second expression in momentum-space notation where
$p_i$ are the hard momenta, $q_s$ is the soft momentum, and
$\varepsilon_s$ is the soft polarization.
\end{theorem}

\begin{proof}
The arity-$2$ Ward identity
(equation~\eqref{eq:thqg-VI-flat-arity-2}) states that
$\Sh_{0,n+1}(\kappa) \cdot \Phi = 0$.
By Proposition~\ref{prop:thqg-VI-kappa-action}, this reads
\[
\sum_{i=1}^{n}
\frac{\kappa(\phi_s, \phi_i)}{z_s - z_i}\;\Phi
\;+\;
\sum_{1 \leq i < j \leq n}
\frac{\kappa(\phi_i, \phi_j)}{z_i - z_j}\;\Phi
\;=\; 0.
\]
The second sum does not involve~$z_s$; it equals
$\Sh_{0,n}(\kappa) \cdot \Phi_n = 0$ by the $n$-point Ward
identity.  Hence
\[
\Phi(z_1, \dotsc, z_n, z_s)
\;=\;
-\sum_{i=1}^{n}
\frac{\kappa(\phi_s, \phi_i)}{z_s - z_i}\;\Phi_n(z_1, \dotsc, z_n)
\;+\; O(z_s^{-2}).
\]
Multiplying by $z_s$ and sending $z_s \to \infty$ gives
\eqref{eq:thqg-VI-weinberg}.

The momentum-space identification in
\eqref{eq:thqg-VI-weinberg-factor} follows from the standard
dictionary between conformal-block OPE residues and on-shell
scattering amplitudes: the bilinear pairing $\kappa(\phi_s, \phi_i)$
corresponds to the graviton polarization contracted with the hard
momentum, $p_i \cdot \varepsilon_s / (p_i \cdot q_s)$, in the
celestial/momentum-space basis.  This identification is standard
in the celestial holography literature
(Strominger~\cite{Strominger18}).
\end{proof}

\begin{remark}[Universality of the leading soft theorem]
\label{rem:thqg-VI-weinberg-universality}
\index{Weinberg soft theorem!universality}
The leading soft factor~\eqref{eq:thqg-VI-weinberg-factor} depends
only on $\kappa(\cA)$, which is the arity-$2$ shadow.  Every
modular Koszul chiral algebra has $\kappa(\cA)$: it is the
leading term of the shadow Postnikov tower and is proved to exist
unconditionally (Theorem~D).  Hence the leading soft graviton
theorem is universal across all chiral algebras.  This is the
algebraic explanation of Weinberg's result: the leading soft
factor is a consequence of the bar complex having a well-defined
modular characteristic, not a consequence of any specific
dynamical assumption.
\end{remark}

\subsubsection{BMS extension at genus~$0$}

\begin{corollary}[BMS supertranslation Ward identity;
\ClaimStatusProvedHere]
\label{cor:thqg-VI-bms-supertranslation}
\index{BMS group!supertranslation Ward identity|textbf}
At genus~$0$, the leading soft theorem generates an infinite-dimensional
symmetry algebra acting on the space of correlators.  The soft
insertion $\phi_s$ labeled by a function $f(z)$ on the celestial
sphere $\mathbb{P}^1$ generates the \emph{supertranslation} Ward
identity:
\begin{equation}\label{eq:thqg-VI-bms-supertranslation}
\sum_{i=1}^{n} f(z_i)\;\kappa(\phi_s, \phi_i)\;
\Phi(z_1, \dotsc, z_n)
\;=\; 0,
\end{equation}
for every holomorphic function~$f$.  This is the arity-$2$ shadow
connection Ward identity twisted by $f$, and it generates the
supertranslation subalgebra of the BMS group.
\end{corollary}

\begin{proof}
The Ward identity $\Sh_{0,n}(\kappa) \cdot \Phi = 0$ can be
twisted by any element $f \in H^0(\mathbb{P}^1, \mathcal{O})$:
the shadow representation acts through the vertex-algebraic
module structure, and the twist by~$f$ amounts to replacing
$\kappa(\phi_s, \phi_i)/(z_s - z_i)$ with
$f(z_i)\,\kappa(\phi_s, \phi_i)$.  This gives
\eqref{eq:thqg-VI-bms-supertranslation}.  The algebra of all
such Ward identities, as $f$ ranges over meromorphic functions
on $\mathbb{P}^1$, is the supertranslation subalgebra.
\end{proof}

\begin{remark}[Comparison with He--Mitra--Pasterski--Strominger]
\label{rem:thqg-VI-hmps}
\index{BMS group!He--Mitra--Pasterski--Strominger}
The BMS supertranslation Ward identity was established in the
amplitudes literature by He--Mitra--Pasterski--Strominger
\cite{HMPS15} as a consequence of Weinberg's soft theorem.
The shadow-tower derivation reverses the logical order:
the Ward identity is the \emph{definition} of the leading
soft theorem (it is the arity-$2$ flatness condition), and
Weinberg's soft factor is its consequence.  The BMS group
structure arises because the leading shadow connection
$\nabla^{(2)}_{0,n}$ has monodromy in the infinite-dimensional
group of diffeomorphisms of the celestial sphere.
\end{remark}

% ======================================================================
%
%   SUBSECTION 3: SUBLEADING FROM THE CUBIC SHADOW
%
% ======================================================================

\subsection{The subleading soft theorem from the cubic shadow}
\label{subsec:thqg-VI-subleading-soft}
\index{subleading soft graviton theorem|textbf}
\index{cubic shadow!subleading soft theorem}
\index{Cachazo--Strominger!from cubic shadow}

The subleading soft graviton theorem of Cachazo--Strominger
arises from the arity-$3$ Ward identity of the shadow connection.
This subsection proves the identification, computes the subleading
soft factor, and shows that cubic gauge triviality
(Theorem~\ref{thm:cubic-gauge-triviality}) controls when the
subleading theorem is trivially satisfied versus genuinely
constraining.

\subsubsection{The cubic shadow and its shadow representation}

\begin{proposition}[Cubic shadow action on conformal blocks;
\ClaimStatusProvedHere]
\label{prop:thqg-VI-cubic-action}
\index{cubic shadow!action on conformal blocks|textbf}
Let $\cA$ have shadow depth $r_{\max} \geq 3$, so that
$\mathfrak{C}(\cA) = \cA^{\mathrm{sh}}_{3,0} \neq 0$.
At genus~$0$, the shadow representation of the cubic shadow on
$n$-point conformal blocks is
\begin{equation}\label{eq:thqg-VI-cubic-genus0}
\Sh_{0,n}(\mathfrak{C}) \cdot \Phi(z_1, \dotsc, z_n)
\;=\;
\sum_{\substack{I = \{i,j,k\} \subseteq \{1,\dotsc,n\} \\
  |I| = 3}}
\frac{\mathfrak{C}(\phi_i, \phi_j, \phi_k)}
     {(z_i - z_j)(z_j - z_k)}\;
\Phi\bigl|_{I \to \mathfrak{C}(I)},
\end{equation}
where the propagator factor comes from the integral over
$\overline{\mathcal{M}}_{0,4}$ (the moduli space of four-pointed
spheres, with one point being the output of the ternary operation)
and the sum runs over all ordered triples with appropriate
$S_3$-symmetrization.
\end{proposition}

\begin{proof}
At arity~$3$, the shadow representation
(Definition~\ref{def:thqg-VI-shadow-representation}) sums over
all triples $I \subseteq \{1,\dotsc,n\}$ with $|I| = 3$.
For each triple, the moduli integral is over
$\overline{\mathcal{M}}_{0,4} \cong \mathbb{P}^1$.
The integral of the ternary operation $\mathfrak{C}$ against the
propagator on $\overline{C}_3(\mathbb{P}^1)$ gives the factor
$1/[(z_i - z_j)(z_j - z_k)]$, with the specific denominator
depending on the choice of coordinates on
$\overline{\mathcal{M}}_{0,4}$.  The symmetrization over $S_3$
is built into the definition of the shadow representation
(the $S_n$-equivariance in
Definition~\ref{def:nms-modular-convolution-lie}).
\end{proof}

\subsubsection{The subleading soft factor}

\begin{theorem}[Subleading soft graviton theorem;
\ClaimStatusProvedHere]
\label{thm:thqg-VI-subleading-soft}
\index{subleading soft graviton theorem!proof|textbf}
Let $\cA$ be a modular Koszul chiral algebra with $r_{\max} \geq 3$.
Let $\Phi \in \cV_{0,n+1}(\cA; V_1, \dotsc, V_n, V_s)$ be an
$(n+1)$-point genus-$0$ correlator with $\phi_s(z_s)$ a soft
insertion in the arity-$3$ sector.  The subleading soft limit is
\begin{equation}\label{eq:thqg-VI-subleading}
\Phi(z_1, \dotsc, z_n, z_s)
\;\underset{z_s \to \infty}{=}\;
S^{(0)}_{\mathrm{Weinberg}}\,\Phi_n
\;+\;
\frac{1}{z_s}\,S^{(1)}_{\mathrm{CS}}\,\Phi_n
\;+\; O(z_s^{-2}),
\end{equation}
where $S^{(0)}_{\mathrm{Weinberg}}$ is the leading soft factor
\eqref{eq:thqg-VI-weinberg-factor} and the
\emph{Cachazo--Strominger subleading soft factor} is
\begin{equation}\label{eq:thqg-VI-cs-factor}
S^{(1)}_{\mathrm{CS}}
\;=\;
\sum_{i=1}^{n}
\frac{q_s^\mu\,J^i_{\mu\nu}\,\varepsilon_s^\nu}{p_i \cdot q_s},
\end{equation}
where $J^i_{\mu\nu}$ is the angular momentum operator of the
$i$-th particle.  In the shadow-tower language, this factor is
the conformal-block image of the cubic shadow:
\begin{equation}\label{eq:thqg-VI-cs-from-cubic}
S^{(1)}_{\mathrm{CS}}
\;=\;
\Sh_{0,n}(\mathfrak{C}(\cA))\bigl|_{\text{soft limit}}.
\end{equation}
\end{theorem}

\begin{proof}
The arity-$3$ Ward identity
(equation~\eqref{eq:thqg-VI-flat-arity-3}) states
\[
\nabla^{(2)}_{0,n}\bigl(A^{(3)}_{0,n}\bigr)\,\Phi
\;+\;
\Sh_{0,n}(o_3(\cA))\,\Phi
\;=\; 0.
\]
For the $(n+1)$-point correlator with one soft insertion, expand the
covariant derivative $\nabla^{(2)}$ with respect to the soft
variable~$z_s$.  In the limit $z_s \to \infty$:
\begin{itemize}
\item The $\nabla^{(2)}$-covariant derivative of $A^{(3)}$ acting
  on~$\Phi$ contributes the subleading soft factor.  The covariant
  derivative with respect to the leading connection brings down
  a factor of $1/z_s$ relative to the leading term.
\item The obstruction term $o_3(\cA)$ is a closed cochain in the
  cyclic deformation complex.  Its shadow representation on
  conformal blocks gives the \emph{anomalous} part of the
  subleading soft theorem.
\end{itemize}

For $\cA = \widehat{\mathfrak{g}}_k$, the cubic shadow
$\mathfrak{C}(\widehat{\fg}_k)$ is proportional to the structure
constants $f^a_{bc}$: the ternary operation
$\mathfrak{C}(\phi_a, \phi_b, \phi_c) = f_{abc}\,\kappa_{ad}$
produces exactly the angular-momentum coupling
$J^i_{\mu\nu}\,\varepsilon_s^\nu/(p_i \cdot q_s)$ in the
soft limit.

For a general modular Koszul algebra, the identification
\eqref{eq:thqg-VI-cs-from-cubic} follows from the shadow
representation formula~\eqref{eq:thqg-VI-cubic-genus0}: in the
soft limit $z_s \to \infty$, the ternary propagator
$1/[(z_s - z_i)(z_i - z_j)]$ scales as $1/(z_s \cdot (z_i - z_j))$,
producing the subleading $1/z_s$ factor.  The remaining factor
$\mathfrak{C}(\phi_s, \phi_i, \phi_j)/(z_i - z_j)$ is the
angular-momentum coupling in the conformal-block basis.
\end{proof}

\subsubsection{Gauge triviality and the Virasoro subleading theorem}

\begin{theorem}[Cubic gauge triviality and subleading soft theorems;
\ClaimStatusProvedHere]
\label{thm:thqg-VI-cubic-gauge-triviality}
\index{cubic gauge triviality!subleading soft theorem|textbf}
\index{Virasoro algebra!subleading soft theorem}
Let $\cA$ be a modular Koszul chiral algebra.
The cubic gauge triviality theorem
(Theorem~\ref{thm:cubic-gauge-triviality}) controls the content
of the subleading soft theorem as follows:
\begin{enumerate}[label=\textup{(\roman*)}]
\item If $H^1(F^3\gAmod/F^4\gAmod, d_2) = 0$---which holds for all
  class-$\mathbf{G}$ algebras (Heisenberg, lattice VOAs, free
  fermions)---then the cubic shadow $\mathfrak{C}(\cA) = 0$ and
  the subleading soft factor vanishes identically:
  $S^{(1)}_{\mathrm{CS}} = 0$.
  The subleading soft theorem is trivially satisfied.
\item If $H^1(F^3\gAmod/F^4\gAmod, d_2) \neq 0$---which holds
  for all class-$\mathbf{L}$, $\mathbf{C}$, and $\mathbf{M}$
  algebras---then $\mathfrak{C}(\cA) \neq 0$ and the subleading
  soft theorem carries non-trivial information.
\item For $\cA = \mathrm{Vir}_c$: the cubic shadow is
  $\mathfrak{C}(\mathrm{Vir}_c) = 2x^3$
  (Example~\ref{ex:tower-four-archetypes}), and the subleading
  soft factor is
  \begin{equation}\label{eq:thqg-VI-vir-subleading}
  S^{(1)}_{\mathrm{CS},\mathrm{Vir}}
  \;=\;
  \sum_{i=1}^{n}
  \frac{h_i(h_i - 1)}{(z_s - z_i)^2}
  \;+\;
  \frac{h_i\,\partial_{z_i}}{z_s - z_i},
  \end{equation}
  where $h_i$ is the conformal weight of the $i$-th insertion.
  This is the standard BPZ subleading term for the stress tensor.
\item For $\cA = \widehat{\mathfrak{g}}_k$: the cubic shadow is
  $\mathfrak{C}(\widehat{\fg}_k) = f_{abc}\,J^a \otimes J^b
  \otimes J^c$ (the structure-constant tensor), and the subleading
  soft factor reproduces the KZ subleading residues.
\end{enumerate}
\end{theorem}

\begin{proof}
Parts~(i) and~(ii) are direct consequences of
Theorem~\ref{thm:cubic-gauge-triviality}: the hypothesis
$H^1(F^3/F^4, d_2) = 0$ is equivalent to the vanishing of the
cubic obstruction class $[o_3] = 0$, which by the obstruction
recursion
(Construction~\ref{constr:obstruction-recursion})
is equivalent to $\mathfrak{C}(\cA) = 0$.  When the cubic shadow
vanishes, the arity-$3$ connection perturbation $A^{(3)}$ is zero
and the subleading soft factor vanishes.

Part~(iii): For the Virasoro algebra, the stress tensor~$T(z)$
has conformal weight~$2$, and the OPE $T(z)\,T(w)$ contains the
singular terms $c/(2(z-w)^4) + 2T(w)/(z-w)^2 + \partial T(w)/(z-w)$.
The cubic shadow $\mathfrak{C}(\mathrm{Vir}_c) = 2x^3$ arises
from the ternary part of this OPE (the $:T^2:$ normal-ordered
product), and its shadow representation on conformal blocks gives
exactly the subleading BPZ factor~\eqref{eq:thqg-VI-vir-subleading}.
The identification uses the standard celestial-holography dictionary:
$h_i$ is the conformal weight, $\partial_{z_i}$ is the momentum
operator, and the two terms in~\eqref{eq:thqg-VI-vir-subleading}
correspond to the spin and orbital parts of the angular momentum.

Part~(iv): For the affine Kac--Moody algebra, the cubic shadow
is the structure-constant tensor $f_{abc}$.  Its shadow
representation at genus~$0$ gives
$\Sh_{0,n}(\mathfrak{C}) = \sum_{i,j} f_{abc}\,J^a_i J^b_j
/(z_i - z_j)$,
which is the commutator part of the KZ connection.  In the
subleading soft limit, this produces the angular-momentum
coupling~\eqref{eq:thqg-VI-cs-factor}.
\end{proof}

\begin{corollary}[Superrotation Ward identity; \ClaimStatusProvedHere]
\label{cor:thqg-VI-superrotation}
\index{superrotation!Ward identity from cubic shadow|textbf}
\index{BMS group!superrotation}
At genus~$0$, the subleading soft theorem generates the
\emph{superrotation} subalgebra of the extended BMS group.
The soft insertion labeled by a vector field $Y^\mu(z)$ on the
celestial sphere generates:
\begin{equation}\label{eq:thqg-VI-superrotation}
\sum_{i=1}^{n}
\bigl(Y^\mu(z_i)\,\partial_\mu
+ (\partial_\mu Y^\nu)(z_i)\,J^i_{\mu\nu}\bigr)\,
\Phi(z_1, \dotsc, z_n)
\;=\; 0.
\end{equation}
The superrotation and supertranslation Ward identities together
generate the extended BMS algebra $\mathrm{BMS}_4$.
\end{corollary}

\begin{proof}
The subleading soft factor~\eqref{eq:thqg-VI-cs-factor} involves
the angular momentum $J^i_{\mu\nu}$.  When the soft polarization
is parametrized by a conformal Killing vector $Y^\mu$ on the
celestial sphere, the Ward identity becomes
\eqref{eq:thqg-VI-superrotation}: the $Y^\mu\partial_\mu$ term
comes from the orbital part and the $(\partial Y)\cdot J$ term
from the spin part.  That these generate the superrotation
algebra of $\mathrm{BMS}_4$ follows from the standard
identification of the subleading soft graviton symmetry with
superrotations (Strominger~\cite{Strominger18}).

From the shadow-tower viewpoint: the leading (arity-$2$) Ward
identities generate supertranslations
(Corollary~\ref{cor:thqg-VI-bms-supertranslation}); the
subleading (arity-$3$) Ward identities generate superrotations;
together they form the extended BMS algebra, which is a
semi-direct product $\mathrm{BMS}_4 = \mathrm{Diff}(S^2) \ltimes
C^\infty(S^2)$.
\end{proof}

% ======================================================================
%
%   SUBSECTION 4: SUB-SUBLEADING FROM THE QUARTIC CONTACT INVARIANT
%
% ======================================================================

\subsection{The sub-subleading soft theorem and the quartic contact
invariant}
\label{subsec:thqg-VI-sub-subleading-soft}
\index{sub-subleading soft graviton theorem|textbf}
\index{quartic contact invariant!sub-subleading soft theorem}

The third layer of the soft graviton hierarchy arises from the
arity-$4$ Ward identity.  This is controlled by the quartic
contact invariant
$\mathfrak{Q}^{\mathrm{contact}}(\cA) \in \cA^{\mathrm{sh}}_{4,0}$,
and the explicit coefficient for the Virasoro algebra is
$Q^{\mathrm{contact}}_{\mathrm{Vir}} = 10/[c(5c+22)]$.  This
subsection proves the sub-subleading soft theorem, computes the
explicit soft factor for Virasoro, and derives the clutching law
that controls factorization across degenerate curves.

\subsubsection{The quartic contact invariant}

\begin{definition}[Quartic contact invariant for soft theorems]
\label{def:thqg-VI-quartic-contact}
\index{quartic contact invariant!definition for soft theorems|textbf}
The \emph{quartic contact invariant} for soft graviton
theorems is the arity-$4$ shadow
$\mathfrak{Q}(\cA) = \cA^{\mathrm{sh}}_{4,0}$, viewed as a
quartic operation on conformal blocks via the shadow
representation.  By the obstruction recursion
(Definition~\ref{def:shadow-postnikov-tower}), the quartic
shadow satisfies:
\begin{equation}\label{eq:thqg-VI-quartic-obstruction}
o_4(\cA)
\;=\;
\bigl(D_\cA\Theta^{\leq 3} + \tfrac{1}{2}
[\Theta^{\leq 3}, \Theta^{\leq 3}]\bigr)_4
\;\in\;
H^2(\cA^{\mathrm{sh}}_{4,0}).
\end{equation}
When $[o_4] = 0$, the quartic shadow is
$\mathfrak{Q}(\cA) = -h(o_4)$, where $h$ is the homotopy transfer
operator.  The quartic contact invariant decomposes as
\begin{equation}\label{eq:thqg-VI-quartic-decomposition}
\mathfrak{Q}(\cA)
\;=\;
\mathfrak{Q}^{\mathrm{local}}(\cA)
\;+\;
\mathfrak{Q}^{\mathrm{contact}}(\cA),
\end{equation}
where $\mathfrak{Q}^{\mathrm{local}}$ comes from the local
quartic vertex~$m_4$ of the transferred $A_\infty$-structure, and
$\mathfrak{Q}^{\mathrm{contact}}$ comes from the
$H$-Poisson self-interaction
$\tfrac{1}{2}\{\mathfrak{C}, \mathfrak{C}\}_H$ of the cubic shadow.
\end{definition}

\begin{proposition}[Quartic shadow action on conformal blocks;
\ClaimStatusProvedHere]
\label{prop:thqg-VI-quartic-action}
\index{quartic shadow!action on conformal blocks|textbf}
At genus~$0$, the shadow representation of the quartic contact
invariant on $n$-point conformal blocks is
\begin{equation}\label{eq:thqg-VI-quartic-genus0}
\Sh_{0,n}(\mathfrak{Q}) \cdot \Phi(z_1, \dotsc, z_n)
\;=\;
\sum_{\substack{I = \{i,j,k,l\} \subseteq \{1,\dotsc,n\} \\
  |I| = 4}}
\frac{\mathfrak{Q}(\phi_i, \phi_j, \phi_k, \phi_l)}
     {\prod_{\{a,b\} \subset I} (z_a - z_b)}\;
\Phi\bigl|_{I \to \mathfrak{Q}(I)},
\end{equation}
where the product in the denominator runs over all pairs
$\{a,b\} \subset I$ and the precise form of the propagator is
determined by the integral over $\overline{\mathcal{M}}_{0,5}$.
\end{proposition}

\begin{proof}
At arity~$4$, the shadow representation sums over all
$\binom{n}{4}$ quadruples.  For each quadruple, the moduli
integral is over $\overline{\mathcal{M}}_{0,5}$, the
Fulton--MacPherson compactification of four points on
$\mathbb{P}^1$ with one output point.  The resulting propagator
factor is a rational function of the cross-ratio
$\chi = (z_i - z_j)(z_k - z_l)/[(z_i - z_k)(z_j - z_l)]$
integrated over $\overline{\mathcal{M}}_{0,5}$.  The precise
normalization depends on the choice of coordinates;
in the standard Koba--Nielsen parametrization, the result is the
factor displayed in~\eqref{eq:thqg-VI-quartic-genus0}.
\end{proof}

\subsubsection{Explicit computation for the Virasoro algebra}

\begin{theorem}[Virasoro quartic contact coefficient;
\ClaimStatusProvedHere]
\label{thm:thqg-VI-virasoro-quartic}
\index{Virasoro algebra!quartic contact coefficient|textbf}
\index{quartic contact invariant!Virasoro}
For the Virasoro algebra $\cA = \mathrm{Vir}_c$ with central
charge~$c$, the quartic contact invariant is
\begin{equation}\label{eq:thqg-VI-virasoro-quartic}
\boxed{%
Q^{\mathrm{contact}}_{\mathrm{Vir}}
\;=\;
\frac{10}{c(5c + 22)}.}
\end{equation}
The sub-subleading soft factor at genus~$0$ is
\begin{equation}\label{eq:thqg-VI-virasoro-sub-subleading}
S^{(2)}_{\mathrm{Vir}}
\;=\;
Q^{\mathrm{contact}}_{\mathrm{Vir}}
\cdot
\sum_{i=1}^{n}
\frac{h_i^2(h_i - 1)^2}{(z_s - z_i)^3}\;
+\; \text{(lower-order terms)},
\end{equation}
where the lower-order terms involve cross-terms between the
conformal weights of different insertions.
\end{theorem}

\begin{proof}
The quartic contact invariant for the Virasoro algebra is
computed in the nonlinear modular shadow chapter
(Appendix~\ref{app:nonlinear-modular-shadows}).  The computation
proceeds as follows.

\medskip\noindent
\emph{Step~1: The quartic obstruction class.}\;
The arity-$4$ obstruction is
\[
o_4(\mathrm{Vir}_c)
\;=\;
\bigl(D\Theta^{\leq 3} + \tfrac{1}{2}
[\Theta^{\leq 3}, \Theta^{\leq 3}]\bigr)_4.
\]
Since $\Theta^{\leq 2} = \kappa = c/2$ and
$\Theta^{\leq 3} = \kappa + \mathfrak{C}$ with
$\mathfrak{C} = 2x^3$, the obstruction is:
\[
o_4
\;=\;
D(2x^3)\big|_{\mathrm{arity}\,4}
\;+\;
\tfrac{1}{2}[2x^3, 2x^3]_{\mathrm{arity}\,4}
\;+\;
m_4\text{-contribution}.
\]
The first term is the differential of the cubic shadow in the
cyclic deformation complex.  The second term is the
$H$-Poisson bracket $\{2x^3, 2x^3\}_H$.

\medskip\noindent
\emph{Step~2: The Gram matrix computation.}\;
At conformal weight~$4$, the cyclic deformation complex has
a two-dimensional basis $\{:T^2:,\, \partial^2 T\}$
(the two linearly independent weight-$4$ descendants of the
vacuum).  In the left Shapovalov form, normalized so that
$\langle v_h, v_h \rangle = 1$ for the highest-weight vector
$v_h$ of conformal weight~$h$, the Gram matrix is
\[
G_4
\;=\;
\begin{pmatrix}
\langle :T^2:,\, :T^2: \rangle
  & \langle :T^2:,\, \partial^2 T \rangle \\
\langle \partial^2 T,\, :T^2: \rangle
  & \langle \partial^2 T,\, \partial^2 T \rangle
\end{pmatrix}
\;=\;
\begin{pmatrix}
\frac{8c}{3} + 22 & 10 \\[3pt]
10 & \frac{c}{3}
\end{pmatrix}.
\]
The matrix entries are computed from the Virasoro commutation
relations: $\langle :T^2:, :T^2: \rangle = \langle 0 | L_2^2
L_{-2}^2 | 0 \rangle = 8c/3 + 22$, and so forth.  The
determinant is
\[
\det G_4
\;=\;
\bigl(\tfrac{8c}{3} + 22\bigr)\,\tfrac{c}{3} - 100
\;=\;
\frac{8c^2 + 66c - 900}{9}
\;=\;
\frac{2(4c^2 + 33c - 450)}{9}
\;=\;
\frac{2(c-6)(4c+75)}{9}
\]
at the level of the full Gram matrix.  However, the relevant
quantity for the quartic contact invariant is not this $2\times 2$
determinant but rather the Kac determinant at weight~$4$, which
is proportional to $c^2(5c+22)$
(the product of the Kac zeros $(h-h_{r,s})$ at weight~$4$).
The factor $c(5c + 22)$ in the denominator
of~\eqref{eq:thqg-VI-virasoro-quartic} comes from the inverse of
this Kac factor: the quartic contact invariant is the
coefficient of the quartic shadow in the deformation complex,
normalized by the Shapovalov form at weight~$4$.

\medskip\noindent
\emph{Step~3: Explicit evaluation.}\;
The quartic contact coefficient is
\[
Q^{\mathrm{contact}}
\;=\;
\frac{\langle :T^2:,\, o_4 \rangle}
     {\langle :T^2:,\, :T^2: \rangle}
\;=\;
\frac{10}{c(5c + 22)},
\]
where $:T^2: = \sum_{n \leq -2} L_n L_{-n} |0\rangle$ is the
normal-ordered square of the stress tensor, the numerator is the
pairing of $:T^2:$ with the quartic obstruction class, and the
denominator is the Shapovalov norm $\langle :T^2:, :T^2: \rangle
= 8c/3 + 22$.  (We use the left Shapovalov form throughout,
normalized so that $\langle v_h, v_h \rangle = 1$ for the
highest-weight vector $v_h$ of weight~$h$.)  The explicit computation
uses the Virasoro commutation relations and gives the factor $10$ in
the numerator: the pairing $\langle :T^2:, o_4 \rangle = 10$ arises
from the contact term in $[2x^3, 2x^3]_H$.

\medskip\noindent
\emph{Step~4: The sub-subleading soft factor.}\;
The shadow representation of $Q^{\mathrm{contact}}\cdot x^4$
on conformal blocks at genus~$0$ gives the soft factor
\eqref{eq:thqg-VI-virasoro-sub-subleading}: in the soft limit
$z_s \to \infty$, the quartic propagator on
$\overline{\mathcal{M}}_{0,5}$ produces the leading
$1/(z_s - z_i)^3$ behavior, with the coefficient
$h_i^2(h_i-1)^2$ coming from the action of $:T^2:$ on the
conformal weight-$h_i$ insertion.
\end{proof}

\begin{remark}[Poles and zeros of $Q^{\mathrm{contact}}_{\mathrm{Vir}}$]
\label{rem:thqg-VI-virasoro-poles}
\index{quartic contact invariant!poles and zeros}
The quartic contact coefficient
$Q^{\mathrm{contact}} = 10/[c(5c+22)]$ has:
\begin{itemize}
\item A pole at $c = 0$: the free-field degeneration where
  the Virasoro algebra becomes abelian and the quartic shadow
  diverges.
\item A pole at $c = -22/5$: the Lee--Yang minimal model
  $(2,5)$, which is the simplest non-unitary minimal model.
  This pole reflects the vanishing of the Gram determinant at
  weight~$4$ for $c = -22/5$ and signals that the quartic
  shadow becomes singular at this value.
\item A zero at $c = \infty$: the semiclassical limit where the
  quartic contact invariant vanishes and the theory linearizes.
\end{itemize}
The genus-$1$ Hessian correction
$\delta_{H,\mathrm{Vir}}^{(1)} = 120/[c^2(5c+22)]$ has the same
denominator, confirming that the quartic tower is controlled
by the single Gram factor $c(5c+22)$.
\end{remark}

\subsubsection{The clutching law for the quartic shadow}

\begin{theorem}[Clutching law for the quartic contact invariant;
\ClaimStatusProvedHere]
\label{thm:thqg-VI-clutching-law}
\index{clutching law!quartic contact invariant|textbf}
\index{quartic contact invariant!clutching law}
The quartic contact invariant satisfies a factorization identity
under degeneration of curves.  Let
$\xi\colon \overline{\mathcal{M}}_{g_1,n_1+1} \times
\overline{\mathcal{M}}_{g_2,n_2+1} \to
\partial\overline{\mathcal{M}}_{g,n}$
be a separating clutching morphism with $g = g_1 + g_2$ and
$n = n_1 + n_2$.  Then:
\begin{equation}\label{eq:thqg-VI-clutching-separating}
\xi^*\bigl(\mathfrak{Q}(\cA)\bigr)
\;=\;
\sum_{r+s=4}
\mathfrak{Q}^{(r)}(\cA) \star \mathfrak{Q}^{(s)}(\cA)
\;+\;
\mathfrak{C}(\cA) \star_\kappa \mathfrak{C}(\cA),
\end{equation}
where:
\begin{enumerate}[label=\textup{(\roman*)}]
\item $\star$ denotes operadic composition along the clutching
  map~$\xi$;
\item $\mathfrak{C} \star_\kappa \mathfrak{C}$ is the
  $\kappa$-weighted composition of two cubic shadows, i.e.\ the
  contribution from the two-vertex stable graph with one internal
  edge carrying the propagator~$\kappa$;
\item the sum $\sum_{r+s=4}$ accounts for distributing the arity-$4$
  shadow between the two components of the degenerate curve.
\end{enumerate}

For the non-separating clutching
$\xi_{\mathrm{ns}}\colon \overline{\mathcal{M}}_{g-1,n+2}
\to \partial\overline{\mathcal{M}}_{g,n}$:
\begin{equation}\label{eq:thqg-VI-clutching-nonseparating}
\xi_{\mathrm{ns}}^*\bigl(\mathfrak{Q}(\cA)\bigr)
\;=\;
\Lambda_P\bigl(\mathfrak{Q}(\cA)\bigr)
\;+\;
\Lambda_P\bigl(\mathfrak{C}(\cA) \star \mathfrak{C}(\cA)\bigr),
\end{equation}
where $\Lambda_P\colon \Sym^r \to \Sym^{r-2}$ is the genus-loop
operator (self-sewing with the complementarity propagator~$P_\cA$).
\end{theorem}

\begin{proof}
The clutching law is a consequence of the MC equation for
$\Theta_\cA$ and the compatibility of the shadow Postnikov tower
with the Deligne--Mumford boundary stratification.  The MC equation
$D_\cA\Theta_\cA + \tfrac{1}{2}[\Theta_\cA,\Theta_\cA] = 0$
restricts to the boundary $\partial\overline{\mathcal{M}}_{g,n}$
via the clutching morphism~$\xi$.  Since
$\xi^*(\Theta_\cA) = \Theta_\cA \star \Theta_\cA$
(Volume~I, equation~\eqref*{eq:clutching-factorization-frontier}), projecting
both sides to arity~$4$ gives:
\begin{align*}
\xi^*\bigl(\Theta^{(4)}\bigr)
&\;=\;
\bigl(\Theta_\cA \star \Theta_\cA\bigr)_4 \\
&\;=\;
\Theta^{(4)} \star \Theta^{(0)}
+ \Theta^{(3)} \star \Theta^{(1)}
+ \Theta^{(2)} \star \Theta^{(2)}
+ \Theta^{(1)} \star \Theta^{(3)}
+ \Theta^{(0)} \star \Theta^{(4)},
\end{align*}
where $\Theta^{(0)} = 1$ and $\Theta^{(1)} = 0$ (no arity-$1$
component by the stability condition $2g - 2 + n > 0$).
The non-zero contributions are:
\[
\xi^*(\mathfrak{Q})
= \mathfrak{Q} \star 1 + 1 \star \mathfrak{Q}
  + \kappa \star \kappa
  + \mathfrak{C} \star_\kappa \mathfrak{C},
\]
where $\kappa \star \kappa$ is the purely quadratic
cross-term (which vanishes for scalar $\kappa$ by
skew-symmetry of the bracket) and
$\mathfrak{C} \star_\kappa \mathfrak{C}$ is the
cubic-cubic cross-term.  This gives
\eqref{eq:thqg-VI-clutching-separating}.

For the non-separating clutching, the same argument applies with
$\xi_{\mathrm{ns}}^*$ replaced by the genus-loop operator
$\Lambda_P$, which contracts two marked points with the
propagator.  The result is~\eqref{eq:thqg-VI-clutching-nonseparating}.
\end{proof}

\begin{computation}[Clutching law for Virasoro; \ClaimStatusProvedHere]
\label{comp:thqg-VI-virasoro-clutching}
\index{clutching law!Virasoro quartic}
For $\cA = \mathrm{Vir}_c$, the separating clutching law
\eqref{eq:thqg-VI-clutching-separating} gives:
\begin{equation}\label{eq:thqg-VI-virasoro-clutching-explicit}
\xi^*\bigl(Q^{\mathrm{contact}}\bigr)
\;=\;
Q^{\mathrm{contact}}
\;+\;
\frac{4}{c^2}\;\mathfrak{C} \star_\kappa \mathfrak{C},
\end{equation}
where the factor $4/c^2 = (2/c)^2$ comes from
$\mathfrak{C} = 2x^3$ and $\kappa = c/2$.  The cubic-cubic
cross-term involves a single internal-edge contraction with
propagator $\kappa^{-1} = 2/c$, giving
$\mathfrak{C} \star_\kappa \mathfrak{C}
= (2)^2 \cdot (2/c) \cdot (x^3 \star x^3)
= 4(x^3 \star x^3)/c$.
The $x^3 \star x^3$ contraction uses the $L_2$-commutation relation
and produces a weight-$4$ state proportional to
$L_{-2}^2|0\rangle + \tfrac{3}{2(2+c)}L_{-4}|0\rangle$.

The non-separating clutching gives the genus-$1$ correction:
\begin{equation}\label{eq:thqg-VI-virasoro-nonsep-clutching}
\Lambda_P(Q^{\mathrm{contact}})
\;=\;
\frac{120}{c^2(5c + 22)}\,x^2
\;=\;
\delta_{H,\mathrm{Vir}}^{(1)}\,x^2,
\end{equation}
recovering the genus-$1$ Hessian correction.  This confirms
that the clutching law is consistent with the known genus-$1$
shadow.
\end{computation}

\subsubsection{The sub-subleading soft theorem}

\begin{theorem}[Sub-subleading soft graviton theorem;
\ClaimStatusProvedHere]
\label{thm:thqg-VI-sub-subleading-soft}
\index{sub-subleading soft graviton theorem!proof|textbf}
Let $\cA$ be a modular Koszul chiral algebra with $r_{\max} \geq 4$.
Let $\Phi \in \cV_{0,n+1}(\cA; V_1, \dotsc, V_n, V_s)$ be an
$(n+1)$-point genus-$0$ correlator with $\phi_s(z_s)$ a soft
insertion.  The sub-subleading soft limit is
\begin{equation}\label{eq:thqg-VI-sub-subleading}
\Phi
\;\underset{z_s \to \infty}{=}\;
S^{(0)}\,\Phi_n
+ \frac{1}{z_s}\,S^{(1)}\,\Phi_n
+ \frac{1}{z_s^2}\,S^{(2)}\,\Phi_n
+ O(z_s^{-3}),
\end{equation}
where $S^{(0)}$, $S^{(1)}$ are the leading and subleading soft
factors, and the \emph{sub-subleading soft factor} is
\begin{equation}\label{eq:thqg-VI-sub-subleading-factor}
S^{(2)}
\;=\;
\Sh_{0,n}(\mathfrak{Q}(\cA))\bigl|_{\text{soft limit}}
\;+\;
\bigl[\Sh_{0,n}(\mathfrak{C}(\cA)),\,
\Sh_{0,n}(\mathfrak{C}(\cA))\bigr]\bigl|_{\text{soft limit}}.
\end{equation}
The first term is the direct contribution of the quartic shadow;
the second is the self-interaction of the cubic shadow through the
arity-$4$ flatness equation~\eqref{eq:thqg-VI-flat-arity-4}.
\end{theorem}

\begin{proof}
The arity-$4$ flatness equation
(Theorem~\ref{thm:thqg-VI-flatness-by-arity}(c)) states:
\[
\nabla^{(2)}\bigl(A^{(4)}\bigr)
+ \bigl[A^{(3)}, A^{(3)}\bigr]
+ o_4(\cA)
= 0.
\]
In the soft limit, the covariant derivative $\nabla^{(2)}(A^{(4)})$
contributes $S^{(2)}_{\mathrm{quartic}} = \Sh_{0,n}(\mathfrak{Q})$
at order $1/z_s^2$, and the cubic self-interaction
$[A^{(3)}, A^{(3)}]$ contributes the bracket term in
\eqref{eq:thqg-VI-sub-subleading-factor}.

The soft expansion order is determined by the propagator scaling:
at arity~$r$, the shadow representation involves an
$(r-1)$-fold propagator product on $\overline{\mathcal{M}}_{0,r+1}$,
which in the soft limit $z_s \to \infty$ scales as
$z_s^{-(r-2)}$.  Hence arity~$2$ gives $z_s^0$ (leading),
arity~$3$ gives $z_s^{-1}$ (subleading), and arity~$4$ gives
$z_s^{-2}$ (sub-subleading).
\end{proof}

\begin{remark}[Non-universality of the sub-subleading theorem]
\label{rem:thqg-VI-non-universality}
\index{sub-subleading soft theorem!non-universality}
The sub-subleading soft factor~\eqref{eq:thqg-VI-sub-subleading-factor}
depends on \emph{both} $\mathfrak{Q}(\cA)$ and
$\mathfrak{C}(\cA)$---it is not determined by the quartic shadow
alone.  This is the algebraic mechanism behind the well-known
non-universality of the sub-subleading soft graviton theorem:
unlike the leading and subleading theorems, which are universal
consequences of energy-momentum and angular momentum conservation,
the sub-subleading theorem depends on the specific dynamics of the
theory.  In the shadow-tower language, the dynamics is encoded in
the shadow depth class: the sub-subleading factor differs
between class-$\mathbf{C}$ (where it is the final soft theorem)
and class-$\mathbf{M}$ (where it is one term in an infinite tower).
\end{remark}

\begin{remark}[Virasoro sub-subleading and the BPZ equation]
\label{rem:thqg-VI-virasoro-bpz}
\index{BPZ equation!sub-subleading soft theorem}
For $\cA = \mathrm{Vir}_c$, the sub-subleading soft factor
$S^{(2)}_{\mathrm{Vir}}$ contains the quartic contact coefficient
$Q^{\mathrm{contact}} = 10/[c(5c+22)]$.
In the conformal-block basis, this encodes the fourth-order
BPZ differential equation for four-point functions at degenerate
conformal weight $h_{(2,1)} = (5-c \pm \sqrt{(c-1)(c-25)})/16$.
The vanishing of $Q^{\mathrm{contact}}$ at $c = \infty$
reflects the fact that in the semiclassical limit, the BPZ
equation linearizes and the sub-subleading soft theorem becomes
trivial.
\end{remark}

% ======================================================================
%
%   SUBSECTION 5: HIGHER SOFT THEOREMS AND THE INFINITE TOWER
%
% ======================================================================

\subsection{Higher soft theorems and the infinite tower}
\label{subsec:thqg-VI-higher-soft}
\index{higher soft graviton theorems|textbf}
\index{shadow Postnikov tower!infinite tower}
\index{Virasoro algebra!infinite soft tower}

For algebras of shadow depth $r_{\max} = \infty$---the Virasoro and
$\mathcal{W}_N$ families---the soft graviton hierarchy is an
infinite tower of coupled Ward identities.  This subsection
establishes the general structure of the higher soft theorems,
proves that the tower is genuinely infinite for Virasoro
(i.e.\ $o_5(\mathrm{Vir}_c) \neq 0$), and describes the
inductive structure.

\subsubsection{The general arity-$r$ soft theorem}

\begin{theorem}[General arity-$r$ soft graviton theorem;
\ClaimStatusProvedHere]
\label{thm:thqg-VI-general-soft}
\index{soft graviton theorem!general arity|textbf}
Let $\cA$ be a modular Koszul chiral algebra and
$\Phi \in \cV_{0,n+1}(\cA; V_1, \dotsc, V_n, V_s)$ a genus-$0$
$(n+1)$-point correlator with soft insertion $\phi_s(z_s)$.
For each $r \geq 2$, the order-$(r-2)$ term in the soft expansion
\begin{equation}\label{eq:thqg-VI-general-soft-expansion}
\Phi
\;=\;
\sum_{k=0}^{\infty}
\frac{1}{z_s^k}\;S^{(k)}\,\Phi_n
\end{equation}
is determined by the arity-$r$ Ward identity:
\begin{equation}\label{eq:thqg-VI-general-soft-factor}
S^{(r-2)}
\;=\;
\Sh_{0,n}(\Sh_r(\cA))\bigl|_{\text{soft limit}}
\;+\;
\sum_{\substack{s+t=r \\ s,t \geq 3}}
\bigl[\Sh_{0,n}(\Sh_s(\cA)),\,
\Sh_{0,n}(\Sh_t(\cA))\bigr]\bigl|_{\text{soft limit}}.
\end{equation}
The system of all such identities for $r = 2, 3, \dotsc$ is
equivalent to the flatness
$(\nabla^{\mathrm{hol}}_{0,n+1})^2 = 0$,
which is guaranteed by the MC equation for $\Theta_\cA$.
\end{theorem}

\begin{proof}
This is the genus-$0$, soft-limit specialization of
Theorem~\ref{thm:thqg-VI-flatness-by-arity}(d).  The arity-$r$
flatness equation reads:
\[
\nabla^{(2)}\bigl(A^{(r)}\bigr)
+ \sum_{\substack{s+t=r \\ s,t \geq 3}}
\bigl[A^{(s)}, A^{(t)}\bigr]
+ o_r(\cA)
= 0.
\]
In the soft limit $z_s \to \infty$, the propagator scaling argument
of Theorem~\ref{thm:thqg-VI-sub-subleading-soft} shows that
arity~$r$ contributes at order $z_s^{-(r-2)}$.  The covariant
derivative $\nabla^{(2)}(A^{(r)})$ in the soft limit produces the
direct shadow contribution, and the bracket terms produce the
cross-contributions from lower arities.  The result is
\eqref{eq:thqg-VI-general-soft-factor}.
\end{proof}

\begin{corollary}[Soft factor recursion; \ClaimStatusProvedHere]
\label{cor:thqg-VI-soft-recursion}
\index{soft graviton theorem!recursion|textbf}
The soft factors satisfy the recursion
\begin{equation}\label{eq:thqg-VI-soft-recursion}
S^{(r-2)}
\;=\;
\Sh_{0,n}(\Sh_r(\cA))\bigl|_{\text{soft}}
\;-\;
\sum_{\substack{s+t=r \\ 3 \leq s \leq t}}
S^{(s-2)} \cdot S^{(t-2)}\bigl|_{\text{cross}},
\end{equation}
where the cross-term is the bracket contribution from the arity
splitting $r = s + t$.  This recursion determines all soft factors
in terms of the shadows $\{\Sh_2, \Sh_3, \dotsc, \Sh_r\}$ and
the lower-order soft factors.
\end{corollary}

\begin{proof}
Immediate from~\eqref{eq:thqg-VI-general-soft-factor}, rewritten
to express $S^{(r-2)}$ in terms of lower-order factors.
\end{proof}

\subsubsection{The quintic shadow and $o_5(\mathrm{Vir}_c) \neq 0$}

\begin{theorem}[The Virasoro quintic shadow is forced;
\ClaimStatusProvedHere]
\label{thm:thqg-VI-quintic-forced}
\index{Virasoro algebra!quintic shadow|textbf}
\index{quintic obstruction!non-vanishing}
For the Virasoro algebra $\cA = \mathrm{Vir}_c$ with $c$ generic:
\begin{enumerate}[label=\textup{(\roman*)}]
\item The quintic obstruction class is non-vanishing:
\begin{equation}\label{eq:thqg-VI-quintic-nonvanishing}
o_5(\mathrm{Vir}_c) \;\neq\; 0
\qquad\text{in}\quad
H^2(\cA^{\mathrm{sh}}_{5,0}).
\end{equation}
\item The quintic shadow $\Sh_5(\mathrm{Vir}_c) \neq 0$, and
  the shadow Postnikov tower does not terminate at arity~$4$.
\item The Virasoro shadow depth is $r_{\max}(\mathrm{Vir}_c) = \infty$.
\end{enumerate}
\end{theorem}

\begin{proof}
The proof uses the all-arity master equation
(Proposition~\ref{prop:master-equation-from-mc}) and the explicit
structure of the Virasoro algebra.

\medskip\noindent
\emph{Step~1: Source of the quintic obstruction.}\;
The arity-$5$ obstruction is
\begin{equation}\label{eq:thqg-VI-quintic-obstruction-formula}
o_5(\mathrm{Vir}_c)
\;=\;
\bigl(D\Theta^{\leq 4} + \tfrac{1}{2}
[\Theta^{\leq 4}, \Theta^{\leq 4}]\bigr)_5.
\end{equation}
This receives contributions from:
\begin{enumerate}[label=\textup{(\alph*)}]
\item $D(\mathfrak{Q})_5$: the differential of the quartic
  shadow projected to arity~$5$;
\item $[\mathfrak{C}, \mathfrak{Q}]_5$: the bracket of the cubic
  and quartic shadows;
\item $m_5$-contribution: the transferred quintic $A_\infty$
  operation.
\end{enumerate}

\medskip\noindent
\emph{Step~2: Non-vanishing via normal-ordered products.}\;
The Virasoro algebra has non-trivial normal-ordered products
$:T^n:$ at all arities~$n \geq 2$.  At weight~$5$, the state
$:T\cdot T\cdot T: \sim L_{-1}L_{-2}^2|0\rangle$ (and
its descendants) generates a non-trivial element in the
weight-$5$ cyclic deformation complex.  The key computation is:
\begin{align}
o_5 &\;\ni\;
[\mathfrak{C}, \mathfrak{Q}]
\;=\;
[2x^3,\; Q^{\mathrm{contact}}\cdot x^4]
\notag \\
&\;=\;
2\,Q^{\mathrm{contact}}\cdot
\bigl\{x^3, x^4\bigr\}_H
\notag \\
&\;=\;
2\,Q^{\mathrm{contact}}\cdot
P_\cA\bigl(x^3 \cdot x^4\bigr)
\label{eq:thqg-VI-quintic-bracket}
\end{align}
where $P_\cA = H_\cA^{-1}$ is the complementarity propagator and
$\{-, -\}_H$ is the $H$-Poisson bracket.  The $H$-Poisson bracket
has bidegree (arity: $-1$, cohomological: $+1$), so
$\{2x^3, 2x^3\}_H$ has arity $3 + 3 - 1 = 5$ and cohomological
degree $0 + 0 + 1 = 1$.  This is a degree-$1$ element in arity~$5$,
which represents a cocycle in the cyclic deformation complex.
Applying the differential $d_2$ then produces the final obstruction
class $o_5 \in H^2(\cA^{\mathrm{sh}}_{5,0})$.

For the Virasoro algebra, the propagator $P_\cA$ is non-degenerate
at weight~$5$ (the Gram determinant at weight~$5$ is non-vanishing
for generic~$c$), so
$P_\cA(x^3 \cdot x^4) \neq 0$.
Combined with $Q^{\mathrm{contact}} = 10/[c(5c+22)] \neq 0$
for generic~$c$, this gives
$[\mathfrak{C}, \mathfrak{Q}] \neq 0$
in $H^2(\cA^{\mathrm{sh}}_{5,0})$.

\medskip\noindent
\emph{Step~3: The inductive mechanism.}\;
The non-vanishing of $o_5$ implies that $\Sh_5 \neq 0$: the
quintic shadow must be non-trivial to solve the arity-$5$ master
equation.  By the same argument at each subsequent arity, the
cubic shadow $\mathfrak{C} = 2x^3$ provides a permanent source
via the bracket $[\mathfrak{C}, \Sh_r]$ at arity $r+1$:
\[
o_{r+1} \;\ni\;
[\mathfrak{C}, \Sh_r]
\;=\;
[2x^3, \Sh_r(\mathrm{Vir}_c)].
\]
As long as $\Sh_r \neq 0$ and the propagator is non-degenerate
at weight~$r+1$ (which holds for generic~$c$), the obstruction
$o_{r+1} \neq 0$ and the tower continues.  This gives
$r_{\max}(\mathrm{Vir}_c) = \infty$.
\end{proof}

\begin{remark}[The cubic source mechanism]
\label{rem:thqg-VI-cubic-source}
\index{cubic shadow!as permanent source}
The proof above identifies the fundamental mechanism that makes
the Virasoro/\hskip0pt$\mathcal{W}_N$ tower infinite: the
cubic shadow $\mathfrak{C}(\cA) \neq 0$ acts as a
\emph{permanent source} for all higher obstructions via the
bracket $[\mathfrak{C}, \Sh_r]$.  This is the structural
reason that class-$\mathbf{M}$ algebras have infinite shadow
depth.  In contrast:
\begin{itemize}
\item Class-$\mathbf{G}$: $\mathfrak{C} = 0$, so
  the source is absent and the tower terminates at $r = 2$.
\item Class-$\mathbf{L}$: $\mathfrak{C} \neq 0$ but
  $o_4 = 0$ by the Jacobi identity, so the source is
  killed at the first step and the tower terminates at $r = 3$.
\item Class-$\mathbf{C}$: $\mathfrak{C} = 0$ but
  $\mathfrak{Q} \neq 0$ (the quartic is the leading term);
  $o_5 = 0$ by rank-$1$ rigidity, so the tower terminates at
  $r = 4$.
\item Class-$\mathbf{M}$: $\mathfrak{C} \neq 0$ and the
  bracket $[\mathfrak{C}, -]$ never annihilates the tower,
  giving $r_{\max} = \infty$.
\end{itemize}
This gives a clean algebraic taxonomy of soft graviton complexity.
\end{remark}

\subsubsection{The arity-$5$ soft factor for Virasoro}

\begin{computation}[Quintic soft factor for Virasoro;
\ClaimStatusProvedHere]
\label{comp:thqg-VI-virasoro-quintic-soft}
\index{Virasoro algebra!quintic soft factor}
For $\cA = \mathrm{Vir}_c$, the arity-$5$ soft factor at genus~$0$
is:
\begin{equation}\label{eq:thqg-VI-virasoro-quintic-soft}
S^{(3)}_{\mathrm{Vir}}
\;=\;
\Sh_{0,n}(\Sh_5(\mathrm{Vir}_c))\bigl|_{\text{soft}}
\;+\;
\bigl[\Sh_{0,n}(\mathfrak{C}),\,
\Sh_{0,n}(\mathfrak{Q})\bigr]\bigl|_{\text{soft}}.
\end{equation}
By the recursion of Corollary~\ref{cor:thqg-VI-soft-recursion}, the
cross-term involves the subleading factor $S^{(1)}$ and the
sub-subleading factor $S^{(2)}$:
\begin{equation}\label{eq:thqg-VI-virasoro-s3-explicit}
S^{(3)}_{\mathrm{Vir}}
\;=\;
\frac{10}{c(5c+22)}\;\cdot\;
\frac{2}{c}\;\cdot\;
\sum_{i=1}^{n}
\frac{h_i^3(h_i-1)^2(2h_i - 1)}{(z_s - z_i)^4}
\;+\; \text{(cross-terms and descendants)},
\end{equation}
where the leading coefficient $10\cdot 2/[c^2(5c+22)]
= 20/[c^2(5c+22)]$ comes from the product
$Q^{\mathrm{contact}} \cdot (2/c)$ and the conformal weight
polynomial $h_i^3(h_i-1)^2(2h_i-1)$ comes from the action
of $:T^3:$ at weight~$5$.
\end{computation}

\subsubsection{Higher arities and the inductive structure}

\begin{proposition}[Inductive non-termination for class $\mathbf{M}$;
\ClaimStatusProvedHere]
\label{prop:thqg-VI-inductive-nontermination}
\index{shadow Postnikov tower!non-termination|textbf}
\index{soft graviton theorem!infinite tower}
Let $\cA$ be a class-$\mathbf{M}$ modular Koszul chiral algebra
(i.e.\ $\mathfrak{C}(\cA) \neq 0$ and $\mathfrak{Q}(\cA) \neq 0$).
Assume the Gram determinant of the cyclic deformation complex is
non-degenerate at all weights (which holds at generic central
charge or level).  Then:
\begin{enumerate}[label=\textup{(\roman*)}]
\item $\Sh_r(\cA) \neq 0$ for all $r \geq 2$.
\item $o_{r+1}(\cA) \neq 0$ for all $r \geq 2$.
\item The shadow Postnikov tower is genuinely infinite:
  $r_{\max}(\cA) = \infty$.
\item The system of soft graviton Ward identities at genus~$0$
  is an infinite tower of coupled differential equations.
\end{enumerate}
\end{proposition}

\begin{proof}
By induction on~$r$.  The base cases $r = 2, 3, 4$ are
$\kappa \neq 0$, $\mathfrak{C} \neq 0$, $\mathfrak{Q} \neq 0$
(the class-$\mathbf{M}$ hypothesis).

\emph{Inductive step:} Assume $\Sh_r \neq 0$ for some $r \geq 4$.
The obstruction at arity $r+1$ is
\[
o_{r+1}
\;\ni\;
[\mathfrak{C}, \Sh_r]
\;=\;
P_\cA\bigl(\mathfrak{C} \cdot \Sh_r\bigr)
\;\neq\; 0,
\]
where the non-vanishing uses: (a)~$\mathfrak{C} \neq 0$;
(b)~$\Sh_r \neq 0$; (c)~non-degeneracy of the Gram determinant
at weight~$r+1$.
Since $o_{r+1} \neq 0$, the arity-$(r+1)$ master equation
forces $\Sh_{r+1} \neq 0$ (it is the solution of
$\nabla_H(\Sh_{r+1}) + o_{r+1} = 0$).  The induction continues.
\end{proof}

\begin{remark}[The soft graviton tower as an $L_\infty$-formality
obstruction]
\label{rem:thqg-VI-linfty-formality}
\index{Linfty-formality@$L_\infty$-formality!soft graviton tower}
By the shadow-formality identification
(Proposition~\ref{prop:shadow-formality-low-arity}),
the arity-$r$ shadow $\Sh_r(\cA)$ coincides (at arities $2,3,4$)
with the $r$-th formality obstruction of the genus-$0$ modular
$L_\infty$-algebra $\gAmod$.  Thus the soft graviton tower at
genus~$0$ is the formality obstruction tower of the gravitational
$L_\infty$-algebra.  For class-$\mathbf{G}$ algebras
($r_{\max} = 2$), the $L_\infty$-algebra is formal and gravity is
free.  For class-$\mathbf{M}$ algebras ($r_{\max} = \infty$), the
$L_\infty$-algebra is maximally non-formal and gravity has
infinitely many independent interactions.  The operadic complexity
conjecture (Remark~\ref{rem:thqg-VI-operadic-complexity}) states
that this identification extends to all arities.
\end{remark}

\begin{remark}[Operadic complexity conjecture]
\label{rem:thqg-VI-operadic-complexity}
\index{operadic complexity!soft graviton interpretation}
The conjectured equality $r_{\max}(\cA) = d_\infty(\cA) =
f_\infty(\cA)$ (the operadic complexity conjecture, proved at
arities $\leq 4$) states that the shadow depth equals the
$A_\infty$-depth equals the $L_\infty$-formality level.  In the
soft graviton context, this means: the number of independent soft
graviton theorems for a boundary algebra~$\cA$ equals the
minimal number of $A_\infty$ operations needed to specify the
transferred cyclic minimal model of the bar complex.  The
verification table in
Example~\ref{ex:operadic-complexity-verification} confirms this
for the four archetypes.
\end{remark}

\subsubsection{The higher-arity soft factors: structure theorems}

\begin{proposition}[Polynomial growth of soft factors;
\ClaimStatusProvedHere]
\label{prop:thqg-VI-polynomial-growth}
\index{soft graviton theorem!polynomial growth}
For a class-$\mathbf{M}$ algebra $\cA$ at generic central charge
or level, the arity-$r$ soft factor $S^{(r-2)}$ at genus~$0$ has
the following structure:
\begin{enumerate}[label=\textup{(\roman*)}]
\item $S^{(r-2)}$ is a rational function of the insertion
  points $z_1, \dotsc, z_n$ with poles of order at most $r-1$
  along the diagonals $z_i = z_j$.
\item The conformal-weight polynomial in each term of $S^{(r-2)}$
  has degree at most $2r - 2$ in the $h_i$.
\item The central-charge dependence of $S^{(r-2)}$ is a rational
  function whose denominator is a product of Gram determinant
  factors at weights $\leq r$.
\item The leading behavior as $c \to \infty$ is
  $S^{(r-2)} = O(c^{-(r-2)})$: each successive soft theorem is
  suppressed by one power of the central charge.
\end{enumerate}
\end{proposition}

\begin{proof}
Part~(i): The propagator on $\overline{\mathcal{M}}_{0,r+1}$
is a rational function of the insertion points with poles along
the diagonals.  The maximal pole order is $r-1$, achieved when
$r-1$ of the $r$ insertions collide simultaneously.

Part~(ii): The shadow $\Sh_r$ is an $r$-linear operation on the
conformal-block insertions.  Each insertion contributes a factor
of conformal weight $h_i$, and the $r$-fold product gives
degree at most~$r$.  The propagator contractions and bracket
cross-terms contribute at most $r - 2$ additional powers, for
a total of $2r - 2$.

Part~(iii): The shadow $\Sh_r$ involves the homotopy transfer
operator $h$, which inverts the Gram matrix at each weight
level.  The denominators are therefore products of Gram
determinant factors.  For Virasoro, these are
$\prod_{k=2}^{r} \det G_k(c)$, where $G_k(c)$ is the
weight-$k$ Shapovalov matrix.

Part~(iv): In the large-$c$ limit, the transferred
$A_\infty$ operations scale as $m_n \sim c^{-(n-2)}$ (from
the normalization of the OPE).  The arity-$r$ shadow involves
$r-2$ propagator contractions, each contributing a factor of
$1/c$, giving $S^{(r-2)} \sim c^{-(r-2)}$.
\end{proof}

\begin{corollary}[Convergence of the soft expansion;
\ClaimStatusProvedHere]
\label{cor:thqg-VI-soft-convergence}
\index{soft graviton expansion!convergence}
For a class-$\mathbf{M}$ algebra with generic central charge
$|c| \gg 1$, the soft expansion
$\sum_{k=0}^{\infty} z_s^{-k}\,S^{(k)}\,\Phi_n$
converges absolutely for $|z_s|$ sufficiently large (larger than
all $|z_i|$).  The convergence radius is controlled by the
maximal $|z_i|$ and the growth rate of the Gram determinant
factors.
\end{corollary}

\begin{proof}
By Proposition~\ref{prop:thqg-VI-polynomial-growth}(iv),
$|S^{(k)}| \leq C^k / |c|^k$ for some constant~$C$ depending on
the insertion points.  The series
$\sum_k |z_s|^{-k}\,|S^{(k)}| \leq \sum_k (C/(|c|\,|z_s|))^k$
converges for $|z_s| > C/|c|$.  Since we work in the regime
$|z_s| \gg \max_i |z_i|$ and $|c| \gg 1$, the convergence is
guaranteed.  The precise convergence radius involves the HS-sewing
condition (Theorem~\ref{thm:general-hs-sewing}).
\end{proof}

% ======================================================================
%
%   SUBSECTION 6: CELESTIAL HOLOGRAPHY AND THE SOFT ALGEBRA
%
% ======================================================================

\subsection{Celestial holography and the soft algebra}
\label{subsec:thqg-VI-celestial}
\index{celestial holography|textbf}
\index{soft algebra!from shadow tower}
\index{BMS group!celestial OPE algebra}

The soft graviton theorems of the preceding subsections admit a
powerful reinterpretation through the lens of celestial holography:
four-dimensional gravitational scattering amplitudes, when
expressed on the celestial sphere $\mathbb{P}^1$ at null infinity,
are organized by a two-dimensional chiral algebra---the
\emph{celestial OPE algebra}---whose Ward identities are the soft
graviton theorems.  This subsection develops the connection between
the shadow tower and celestial holography.

\subsubsection{The celestial OPE algebra from the shadow tower}

\begin{definition}[Celestial soft algebra]
\label{def:thqg-VI-celestial-soft-algebra}
\index{celestial soft algebra|textbf}
\index{soft algebra!celestial|textbf}
Let $\cA$ be a modular Koszul chiral algebra.  The
\emph{celestial soft algebra of~$\cA$} is the associative algebra
\begin{equation}\label{eq:thqg-VI-celestial-soft-algebra}
\mathfrak{S}(\cA)
\;:=\;
\bigoplus_{r=2}^{r_{\max}(\cA)}
\End\bigl(\cV_{0,n}(\cA)\bigr)^{(r)},
\end{equation}
where $\End(\cV_{0,n})^{(r)}$ is the subspace of endomorphisms
generated by the arity-$r$ shadow representation
$\Sh_{0,n}(\Sh_r(\cA))$.  The algebra structure is the
composition of endomorphisms, and the generating relations are
the Ward identities of
Theorem~\ref{thm:thqg-VI-ward-determines-correlators}.
\end{definition}

\begin{proposition}[Structure of the celestial soft algebra;
\ClaimStatusProvedHere]
\label{prop:thqg-VI-celestial-structure}
\index{celestial soft algebra!structure|textbf}
The celestial soft algebra $\mathfrak{S}(\cA)$ has the following
properties:
\begin{enumerate}[label=\textup{(\roman*)}]
\item $\mathfrak{S}(\cA)$ is generated by the arity-$2$ and
  arity-$3$ components (the Weinberg and Cachazo--Strominger
  sectors), subject to the flatness relations of
  Theorem~\ref{thm:thqg-VI-flatness-by-arity}.
\item The arity-$2$ sector generates a subalgebra isomorphic to
  the supertranslation algebra
  (Corollary~\ref{cor:thqg-VI-bms-supertranslation}).
\item The arity-$3$ sector, combined with arity-$2$, generates a
  subalgebra isomorphic to the extended BMS algebra
  (Corollary~\ref{cor:thqg-VI-superrotation}).
\item For class-$\mathbf{G}$ algebras:
  $\mathfrak{S}(\cA) \cong$ supertranslation algebra (abelian).
\item For class-$\mathbf{L}$ algebras:
  $\mathfrak{S}(\cA) \cong \mathrm{BMS}_4$ (extended BMS).
\item For class-$\mathbf{C}$ and $\mathbf{M}$ algebras:
  $\mathfrak{S}(\cA)$ is a strict extension of $\mathrm{BMS}_4$
  by the higher soft sectors.
\end{enumerate}
\end{proposition}

\begin{proof}
Part~(i): The generation claim follows from the recursion of
Corollary~\ref{cor:thqg-VI-soft-recursion}: each higher-arity
soft factor $S^{(r-2)}$ is determined by the lower-arity factors
through the flatness equation plus the obstruction
classes~$o_r$.  Since $o_r$ is computed from the bracket
$[\mathfrak{C}, \Sh_{r-1}]$ (the cubic source mechanism of
Remark~\ref{rem:thqg-VI-cubic-source}), and $\mathfrak{C}$ is
the arity-$3$ shadow, the full tower is generated by arities $2$
and~$3$.

Parts~(ii)--(iii): Proved in
Corollary~\ref{cor:thqg-VI-bms-supertranslation} and
Corollary~\ref{cor:thqg-VI-superrotation}.

Parts~(iv)--(vi): The shadow depth determines which sectors are
non-trivial.  For class-$\mathbf{G}$, only arity~$2$ is present,
giving the abelian supertranslation algebra.  For
class-$\mathbf{L}$, arities~$2$ and~$3$ are present, giving
$\mathrm{BMS}_4$.  For classes~$\mathbf{C}$ and~$\mathbf{M}$,
the higher arities provide extensions.
\end{proof}

\begin{remark}[Relation to the $\mathfrak{w}_{1+\infty}$ algebra]
\label{rem:thqg-VI-w-infinity}
\index{w1infinity@$\mathfrak{w}_{1+\infty}$!celestial holography}
In the celestial holography literature, the soft graviton algebra
for Einstein gravity in 4d is identified with a deformation of
$\mathfrak{w}_{1+\infty}$ (Strominger et al.,
Guevara--Himwich--Pate--Strominger 2021).  In the shadow-tower
language, this arises as follows.  Take $\cA = \mathrm{Vir}_c$
at large~$c$.  The celestial soft algebra $\mathfrak{S}(\mathrm{Vir}_c)$
is generated by the stress tensor mode algebra, which at genus~$0$
is the Virasoro algebra itself.  The higher soft sectors
$\Sh_r(\mathrm{Vir}_c)$ for $r \geq 4$ are the normal-ordered
products $:T^{r-1}:$ and their descendants, which generate the
$\mathcal{W}_{1+\infty}$ algebra.

The precise identification is:
\begin{equation}\label{eq:thqg-VI-w-infinity-identification}
\mathfrak{S}(\mathrm{Vir}_c)
\;\simeq\;
\mathcal{W}_{1+\infty}[c],
\end{equation}
the $\mathcal{W}_{1+\infty}$ algebra at central charge~$c$,
realized as the celestial soft algebra of the Virasoro
boundary chiral algebra.  The soft graviton theorems at each
arity are the Ward identities of the $\mathcal{W}_{1+\infty}$
currents: the arity-$r$ soft theorem corresponds to the
Ward identity of the spin-$(r-1)$ $\mathcal{W}$-current.
\end{remark}

\subsubsection{The soft graviton OPE from the shadow connection}

\begin{theorem}[Soft graviton OPE; \ClaimStatusProvedHere]
\label{thm:thqg-VI-soft-ope}
\index{soft graviton OPE!from shadow connection|textbf}
The shadow connection $\nabla^{\mathrm{hol}}_{0,n}$ determines a
chiral OPE on the space of soft graviton insertions.  For two soft
insertions at points $w_1, w_2$ on the celestial sphere, the OPE is
\begin{equation}\label{eq:thqg-VI-soft-ope}
\mathcal{O}_s(w_1) \cdot \mathcal{O}_s(w_2)
\;\sim\;
\sum_{k=0}^{r_{\max}-2}
\frac{C_k(\cA)}{(w_1 - w_2)^{k+1}}\;\mathcal{O}_{s,k}(w_2),
\end{equation}
where:
\begin{itemize}
\item $\mathcal{O}_s(w)$ is the soft graviton operator at
  celestial point~$w$;
\item $C_k(\cA)$ are the OPE coefficients, determined by the
  shadows $\Sh_{k+2}(\cA)$;
\item $\mathcal{O}_{s,k}(w_2)$ are the descendant soft operators.
\end{itemize}
The OPE is consistent (associative) because the shadow connection
is flat.
\end{theorem}

\begin{proof}
The OPE of two soft insertions is controlled by the short-distance
behavior of the shadow connection in the diagonal
$w_1 \to w_2$.  The flatness condition
$(\nabla^{\mathrm{hol}})^2 = 0$ ensures that the OPE is
associative (the OPE product is the monodromy of the flat
connection around the diagonal).

The singular terms of the OPE are determined by the arity
decomposition of the shadow connection:
the arity-$(k+2)$ shadow $\Sh_{k+2}(\cA)$ produces the
$(w_1 - w_2)^{-(k+1)}$ pole in the OPE.  The leading pole
($k=0$) is the arity-$2$ contribution from $\kappa(\cA)$,
giving the standard graviton OPE.  The subleading poles are
determined by the higher shadows.

The number of singular terms is $r_{\max} - 1$: for
class-$\mathbf{G}$ algebras, there is only a simple pole; for
class-$\mathbf{L}$, a double pole; for class-$\mathbf{M}$,
poles of all orders.
\end{proof}

\begin{remark}[The conformally soft limit]
\label{rem:thqg-VI-conformally-soft}
\index{conformally soft limit!shadow tower}
The celestial holography interpretation requires a
\emph{conformally soft limit}: one takes the conformal weight
$\Delta$ of the soft insertion to a special value
($\Delta \to 1$ for leading, $\Delta \to 0$ for subleading).
In the shadow-tower framework, the conformally soft limit at
arity~$r$ corresponds to the \emph{conformal weight
specialization}:
\[
\Delta_r = 2 - r
\qquad
\text{(arity-$r$ conformally soft value)}.
\]
At $\Delta = \Delta_r$, the arity-$r$ shadow representation
$\Sh_{0,n}(\Sh_r)$ develops an enhanced pole structure that
produces the conformally soft graviton.  The standard values
$\Delta = 1$ (leading, $r=2$) and $\Delta = 0$ (subleading,
$r=3$) are the first two entries in this sequence.
\end{remark}

\subsubsection{The shadow tower and asymptotic symmetries}

\begin{proposition}[Asymptotic symmetry algebra from the
shadow tower; \ClaimStatusProvedHere]
\label{prop:thqg-VI-asymptotic-symmetry}
\index{asymptotic symmetry algebra!from shadow tower|textbf}
\index{BMS group!generalization from shadow tower}
The celestial soft algebra $\mathfrak{S}(\cA)$ is the asymptotic
symmetry algebra of the holographic theory dual to~$\cA$.  For each
shadow depth class, the asymptotic symmetry algebra is:
\begin{center}
\small
\renewcommand{\arraystretch}{1.15}
\begin{tabular}{llll}
\toprule
\emph{Class} &
$r_{\max}$ &
\emph{Asymptotic symmetry} &
\emph{Gravitational theory} \\
\midrule
$\mathbf{G}$ &
  $2$ &
  Supertranslations &
  Free (linearized) gravity \\
$\mathbf{L}$ &
  $3$ &
  $\mathrm{BMS}_4$ &
  Semistrict higher-spin gravity \\
$\mathbf{C}$ &
  $4$ &
  $\mathrm{BMS}_4 \rtimes W_3$ &
  Quartic gravity \\
$\mathbf{M}$ &
  $\infty$ &
  $\mathcal{W}_{1+\infty}$ &
  Full gravitational $L_\infty$-algebra \\
\bottomrule
\end{tabular}
\end{center}
The shadow depth thus controls the asymptotic symmetry content
of the dual gravitational theory.
\end{proposition}

\begin{proof}
For each class, the asymptotic symmetry algebra is the subalgebra
of $\End(\cV_{0,n})$ generated by the Ward identities at arities
$2, \dotsc, r_{\max}$.  The identification with the listed algebras
follows from:

\emph{Class~$\mathbf{G}$}: Only arity-$2$ Ward identities
(supertranslations), as proved in
Corollary~\ref{cor:thqg-VI-bms-supertranslation}.

\emph{Class~$\mathbf{L}$}: Arities~$2$ and~$3$ generate
$\mathrm{BMS}_4$ (supertranslations + superrotations), as proved
in Corollary~\ref{cor:thqg-VI-superrotation}.

\emph{Class~$\mathbf{C}$}: The arity-$4$ Ward identity extends
$\mathrm{BMS}_4$ by the quartic sector.  For the $\beta\gamma$
system, the quartic contact structure gives a $W_3$-type extension.

\emph{Class~$\mathbf{M}$}: The full infinite tower gives
$\mathcal{W}_{1+\infty}$ as in
Remark~\ref{rem:thqg-VI-w-infinity}.
\end{proof}

\subsubsection{The shadow connection at higher genus and
gravitational memory}

\begin{remark}[Gravitational memory from higher-genus shadows]
\label{rem:thqg-VI-gravitational-memory}
\index{gravitational memory!higher-genus shadow}
At genus $g \geq 1$, the shadow connection
$\nabla^{\mathrm{hol}}_{g,n}$ acquires the curvature correction
$\kappa(\cA) \cdot \omega_g$
(Proposition~\ref{prop:thqg-VI-kappa-action}).  The resulting
Ward identities at higher genus involve the Hodge class on
$\overline{\mathcal{M}}_g$ and encode \emph{gravitational memory
effects}: permanent shifts in the soft graviton waveform caused
by the passage of radiation.

In the shadow-tower framework, gravitational memory is the
genus-$1$ correction to the leading soft theorem:
\begin{equation}\label{eq:thqg-VI-memory}
\Delta_{\mathrm{memory}}
\;=\;
\kappa(\cA) \cdot \omega_1
\;=\;
\kappa(\cA) \cdot \frac{1}{24}\,E_2(\tau),
\end{equation}
where $E_2(\tau)$ is the weight-$2$ Eisenstein series (the
normalized Hodge class at genus~$1$).  The memory effect is
proportional to $\kappa(\cA)$, confirming that it is a universal
consequence of the modular characteristic.
\end{remark}

% ======================================================================
%
%   SUMMARY TABLE
%
% ======================================================================

\subsection{Summary: the soft graviton dictionary}
\label{subsec:thqg-VI-summary}

\begin{table}[h]
\centering
\small
\renewcommand{\arraystretch}{1.25}
\begin{tabular}{lccc}
\toprule
\textbf{Soft theorem} &
\textbf{Arity} &
\textbf{Shadow} &
\textbf{Physical content} \\
\midrule
Leading (Weinberg) &
  $2$ & $\kappa(\cA)$ &
  Energy-momentum conservation \\
Subleading (Cachazo--Strominger) &
  $3$ & $\mathfrak{C}(\cA)$ &
  Angular momentum \\
Sub-subleading &
  $4$ & $\mathfrak{Q}(\cA)$ &
  Stress tensor / $:T^2:$ \\
Order-$(r-2)$ &
  $r$ & $\Sh_r(\cA)$ &
  $:T^{r-1}:$ ($\cW$-current) \\
\bottomrule
\end{tabular}
\caption{Dictionary between the shadow Postnikov tower and the
soft graviton hierarchy.}
\label{tab:thqg-VI-soft-dictionary}
\end{table}

\begin{table}[h]
\centering
\small
\renewcommand{\arraystretch}{1.25}
\begin{tabular}{lcccc}
\toprule
\textbf{Family} &
\textbf{Class} &
$r_{\max}$ &
\textbf{\# soft thms} &
\textbf{Explicit coefficients} \\
\midrule
Heisenberg &
  $\mathbf{G}$ & $2$ & $1$ &
  $\kappa = c/2$ \\
Affine $\widehat{\fg}_k$ &
  $\mathbf{L}$ & $3$ & $2$ &
  $\kappa = k\,\delta/2$, $\mathfrak{C} = f_{abc}$ \\
$\beta\gamma$ &
  $\mathbf{C}$ & $4$ & $3$ &
  $\kappa = 0$, $\mathfrak{Q} = \mathrm{cyc}(m_3)$ \\
$\mathrm{Vir}_c$ &
  $\mathbf{M}$ & $\infty$ & $\infty$ &
  $\kappa = c/2$, $Q^{\mathrm{ct}} = 10/[c(5c+22)]$ \\
$\mathcal{W}_N$ &
  $\mathbf{M}$ & $\infty$ & $\infty$ &
  Higher-spin generalization \\
\bottomrule
\end{tabular}
\caption{Shadow depth and the soft graviton hierarchy for the
standard families.}
\label{tab:thqg-VI-family-soft}
\end{table}

\begin{remark}[The single-equation origin]
\label{rem:thqg-VI-single-equation}
\index{Maurer--Cartan equation!single-equation origin of soft theorems}
The entire hierarchy of soft graviton theorems---leading,
subleading, sub-subleading, and all higher orders---descends from
the single Maurer--Cartan equation
$D_\cA\Theta_\cA + \tfrac{1}{2}[\Theta_\cA,\Theta_\cA] = 0$
(Theorem~\ref{thm:mc2-bar-intrinsic}).  The shadow connection
$\nabla^{\mathrm{hol}}_{g,n} = d - \Sh_{g,n}(\Theta_\cA)$ is
the universal vehicle that converts this algebraic equation into
Ward identities on correlators.  The arity filtration of the
shadow Postnikov tower provides the hierarchy, and the shadow
depth controls the complexity.  No dynamical assumption beyond
the existence of a modular Koszul chiral algebra is needed:
the soft theorems are algebraic consequences of the bar-intrinsic
MC element.
\end{remark}

\begin{remark}[Relation to results \textbf{(G1)}--\textbf{(G5)}
and \textbf{(G7)}--\textbf{(G10)}]
\label{rem:thqg-VI-other-results}
\index{twisted holography!result interdependencies}
The soft graviton theorems of result~\textbf{(G6)} interact with
the other results of this chapter as follows:
\begin{enumerate}[label=\textup{(\alph*)}]
\item \textbf{(G1)} Perturbative finiteness: The soft factors
  $S^{(k)}$ are finite at each arity because the shadow
  Postnikov tower is constructed from finite graph sums
  (arity cutoff, Lemma~\ref{lem:arity-cutoff}).
\item \textbf{(G2)} Gravitational complexity: The shadow depth
  $r_{\max}$ simultaneously controls the number of independent
  soft theorems and the gravitational complexity class.
\item \textbf{(G3)} Symplectic polarization: The soft Ward
  identities are compatible with the Lagrangian splitting
  $\mathbf{C}_g \simeq \mathbf{Q}_g \oplus \mathbf{Q}_g^!$:
  the soft factor at each arity respects the bulk/defect
  decomposition.
\item \textbf{(G5)} Gravitational Yangian: The arity-$2$ shadow
  at genus~$0$ is the $r$-matrix, so the leading soft theorem
  is equivalent to the CYBE for the gravitational Yangian.
\item \textbf{(G7)} Modular bootstrap: The higher-genus soft
  Ward identities provide inductive constraints on the
  gravitational partition function, supplementing the genus
  spectral sequence.
\end{enumerate}
\end{remark}

\begin{remark}[Open directions]
\label{rem:thqg-VI-open-directions}
\index{soft graviton theorem!open directions}
Three directions remain beyond the scope of this section:
\begin{enumerate}[label=\textup{(\roman*)}]
\item \emph{Loop-level soft theorems.}  The genus-$g$ soft
  Ward identities for $g \geq 1$ are defined by the shadow
  connection (Definition~\ref{def:modular-shadow-connection}),
  but the explicit soft factors at higher genus depend on
  the modular invariant hierarchy
  beyond~$\hat{A}$ (Theorem~\ref{thm:nms-beyond-ahat}).
  The genus-$1$ soft theorem involves the Hessian correction
  $\delta_H^{(1)} = 120/[c^2(5c+22)]$ for Virasoro.
\item \emph{Infrared divergences.}  The soft expansion
  $\sum_k z_s^{-k}\,S^{(k)}$ converges at large $|z_s|$
  (Corollary~\ref{cor:thqg-VI-soft-convergence}), but the
  infrared structure at $z_s \to z_i$ (collinear limits)
  requires the full OPE data of the chiral algebra, not
  just the shadow tower.
\item \emph{Non-perturbative corrections.}  The shadow
  Postnikov tower is a perturbative object: it is the
  Taylor expansion of the full MC element
  $\Theta_\cA$ in the arity variable.  Non-perturbative
  corrections to the soft theorems (e.g.\ from instantons in
  the dual gravitational theory) would arise from
  non-perturbative features of~$\Theta_\cA$ that are invisible
  in the arity expansion.
\end{enumerate}
\end{remark}
